\documentclass{article}

% ============================================================
% NeurIPS 2026 main-track submission (single-column, double-blind)
% During drafting we keep author info visible via the [preprint] option.
% For the actual submission, switch to plain  \usepackage{neurips_2026}.
% For the camera-ready,     switch to         \usepackage[final]{neurips_2026}.
% ============================================================
\PassOptionsToPackage{numbers,compress}{natbib}
\usepackage[preprint]{neurips_2026}

% --- Core packages ---
\usepackage[utf8]{inputenc}
\usepackage[T1]{fontenc}
\usepackage{hyperref}
\usepackage{url}
\usepackage{booktabs}
\usepackage{amsmath,amssymb,amsfonts}
\usepackage{mathtools}
\usepackage{nicefrac}
\usepackage{microtype}
\usepackage{xcolor}
\usepackage{graphicx}
\usepackage{subcaption}
\usepackage{multirow}
\usepackage{array}
\usepackage{caption}
\usepackage{enumitem}

% --- Convenient macros ---
\newcommand{\Wdown}{\mathbf{W}_{\text{down}}}
\newcommand{\ma}{\mathrm{MA}}
\newcommand{\ffd}{d_{\text{ff}}}
\newcommand{\topk}{\mathrm{Top}}

% ============================================================
% Title & authors
% ============================================================
\title{Function Words as Geometric Anchors:\\
An SVD Expansion Mechanism\\
for Massive Activations in Large Language Models}

\author{%
  David S.~Hippocampus\thanks{Use footnote for providing further information
    about author (webpage, alternative address)---\emph{not} for acknowledging
    funding agencies.} \\
  Department of Computer Science\\
  Cranberry-Lemon University\\
  Pittsburgh, PA 15213 \\
  \texttt{hippo@cs.cranberry-lemon.edu} \\
  % examples of more authors
  % \And
  % Coauthor \\
  % Affiliation \\
  % Address \\
  % \texttt{email} \\
  % \AND
  % Coauthor \\
  % Affiliation \\
  % Address \\
  % \texttt{email} \\
  % \And
  % Coauthor \\
  % Affiliation \\
  % Address \\
  % \texttt{email} \\
  % \And
  % Coauthor \\
  % Affiliation \\
  % Address \\
  % \texttt{email} \\
}

\begin{document}
\maketitle

% ============================================================
% Abstract
% ============================================================
\begin{abstract}
Massive activations (MAs) are hidden-state components whose magnitudes are
orders of magnitude larger than typical activations. They reduce the
stability of low-precision inference and complicate quantisation, yet
have so far lacked a geometric and mechanistic explanation. Our work
provides a unified, causally validated account of MA generation across
a wide range of mainstream large language models, spanning ten
architectural families and three activation functions. We combine
spectral analysis of the MLP down-projection with systematic causal
interventions on attention, MLP sub-modules, right-singular matrices, and
single- to multi-layer MLPs, and show that MAs are produced through
the conjunction of low-entropy function-token positions in the input
and a low-rank subspace of the learned down-projection. Our framework
delivers a mechanistic account of MA generation, a closed-form
signed-change prediction under a randomised down-projection
right-singular matrix, mechanism-level classifications that are
orthogonal to normalisation and activation choices, and a system-level
emergent phenomenon in which MAs emerge across multiple layers via
residual-stream propagation. Our code is available at
\url{https://anonymous.4open.science/r/changeHead_massvieAcitve-A8C6}.
\end{abstract}

% ============================================================
% 1. Introduction
% ============================================================
\section{Introduction}
\label{sec:intro}

Transformer language models~\citep{vaswani2017attention,brown2020language,touvron2023llama}
exhibit massive activations (MAs)---hidden-state components orders of
magnitude larger than typical values~\citep{sun2024massive,merrill2022saturated}.
MAs concentrate on a handful of tokens (often whitespace, punctuation, or the
beginning-of-sequence token) and a handful of hidden dimensions, yet they sit
on the critical path of every forward pass. They destabilise low-precision
inference~\citep{dettmers2022gpt3,bondarenko2023quantizable}, interact
with attention sinks~\citep{xiao2024efficient}, complicate long-context
reasoning~\citep{liu2024lost}, and obstruct quantisation, pruning, or
sharding without retraining~\citep{liu2024alberta}.

Despite this operational salience, what causes an MA in a particular
layer and on a particular token is not settled. Existing accounts fall into
two largely descriptive camps. The first, rooted in the analysis of
attention sinks, correlates MAs with function words and short
structural tokens~\citep{xiao2024efficient,cancedda2024spectral,gu2025attention};
the second, rooted in systems work, treats MAs as numerical artefacts to be
contained via precision-aware quantisation, smoothing, or selective
preservation~\citep{dettmers2022gpt3,xiao2023smoothquant,bondarenko2023quantizable,liu2024intactkv}.
Neither yields a mechanistic explanation of why MAs arise only at
certain depths, concentrate on a handful of hidden dimensions, and respond
sharply to small geometric interventions. A unified theory should (i) predict
where in the network MAs emerge, (ii) predict how large they
become, and (iii) establish causality in a way that scales across
architectures.

We study MAs through mechanistic
interpretability~\citep{elhage2021mathematical,conmy2023towards} and
causal intervention~\citep{pearl2009causality,meng2022locating}. Our
analysis centres on the singular-value decomposition of the MLP
down-projection: for each token we quantify how much of its massive
activation is carried by the leading singular directions, and we
complement this with a $V$-matrix ablation that replaces the
right-singular matrix by a random orthogonal matrix while keeping the
spectrum and the rest of the network fixed. This intervention
disentangles the directional content of the down-projection
from its spectral amplification capacity and, crucially, admits
a closed-form prediction in the high-dimensional regime.

Within this geometric framing, function and content tokens play
asymmetric roles. Function words (closed-class high-frequency items
such as ``the'', ``of'', ``a'', together with structural punctuation
and whitespace tokens) recur in highly similar contexts, so their MLP
intermediate states $\mathbf{h}_2$ concentrate on a low-rank
submanifold of hidden space; content words, by virtue of their
semantic and contextual diversity, induce $\mathbf{h}_2$ that span a
substantially higher-dimensional region. When the top right-singular
directions of the down-projection align with the low-rank subspace
occupied by function-token intermediates, input-side low rank
and weight-side low rank couple: spectral amplification and
directional alignment compound on a single hidden channel and produce
an activation orders of magnitude above baseline. Content tokens,
lacking systematic alignment with the same directions, almost never
trigger MAs at the same layer. This upgrades the function-word
correlation reported in prior
work~\citep{sun2024massive,xiao2024efficient} from a statistical
observation to a quantitative geometric mechanism.

Our central claim has two parts. First, the MLP contribution to an MA
obeys the SVD identity (Eq.~\eqref{eq:ma-poly}, which holds for all
$\Wdown$ when summed to $K=r$). The empirical, falsifiable part is
that a low-rank truncation at $K\leq 20$ already captures most of the
MA magnitude on a wide range of architectures.
The MA-channel MLP output decomposes into a weighted sum of
the token's projections onto the right singular directions of the
down-projection, with weights equal to the corresponding singular values
and left-singular coordinates. When the top singular direction dominates
and the MLP intermediate aligns with it, the expansion reduces to a
one-parameter linear relationship that can be fit directly against
observed activations. Models in which no
single direction dominates are handled by a macro variant of
the $V$-ablation and by a layer-support taxonomy whose measured spectral
dominance and directional alignment explain the observed ablation behaviour.
Together, these results
predict where MAs emerge and how large they become, and
they survive ablation testing across normalisation schemes, activation
functions, and architectural families.
Put compactly, our experiments trace a single generation pathway:
function-token triggering, MLP down-projection generation, SVD-mediated
geometric amplification, and residual-stream propagation.

\paragraph{Contributions.}
\begin{enumerate}[leftmargin=*,itemsep=1pt,topsep=2pt]
  \item \textbf{An SVD expansion and closed-form signed-change
  prediction}: from the SVD of the MLP down-projection, each MA-channel
  MLP output decomposes (as an algebraic identity) into a sum of singular-direction
  projections weighted by their singular values and left-singular
  coordinates. Replacing the right-singular matrix $\mathbf{V}$ with a
  random orthogonal matrix then yields a closed-form expression for the
  expected signed change $\mathbb{E}[\Delta_V]$
  (Eq.~\eqref{eq:delta-v}) that closely matches the empirical
  $\Delta_V$.
  \item \textbf{A geometric mechanism for the function-word preference}:
  MAs concentrate almost exclusively on low-entropy function-token
  positions because the MLP intermediate states at those positions
  occupy a low-rank subspace that aligns with the leading
  right-singular directions of the down-projection, supplying a
  quantitative geometric basis for the attention-sink correlations
  reported in prior work.
  \item \textbf{A four-regime cross-architectural taxonomy}
  (\textsc{Concentrated}, \textsc{Few-Source}, \textsc{Dispersed},
  \textsc{Anomaly}) based on layer support, with spectral dominance and
  directional alignment measured as consequences, which classifies every model in
  our pool without recourse to normalisation or activation-function labels.
  \item \textbf{MAs as a residual-stream-level multi-layer emergent
  phenomenon}: with all other MLPs ablated, restoring only the
  trigger layer rarely regenerates the MA and succeeds only in
  architectural exceptions such as parallel blocks or very early
  trigger layers, showing that the trigger layer depends on
  $\mathbf{h}_2$ accumulated from upstream MLPs through the residual
  stream and that single-layer geometry is load-bearing but not
  sufficient.
\end{enumerate}

% ============================================================
% 2. Related work
% ============================================================
\section{Related Work}
\label{sec:related}

\paragraph{Massive activations.}
\citet{sun2024massive} systematically catalogued MAs in modern language
models and their correlation with function words. Subsequent work studies
their impact on post-training quantisation~\citep{dettmers2022gpt3,xiao2023smoothquant,bondarenko2023quantizable,shao2024omniquant},
accelerator error tolerance~\citep{liu2024alberta}, long-context
behaviour~\citep{liu2024lost}, and pivot-token
preservation~\citep{liu2024intactkv}; earlier, \citet{merrill2022saturated}
analysed activation saturation from a formal-linguistic perspective. These
works treat MAs mainly as a systems phenomenon to be controlled, whereas we
ask why MAs arise and test the causal generation pathway.

\paragraph{Function words, syntactic probes, and attention sinks.}
Probing work shows that function words carry disproportionate syntactic
signal~\citep{hewitt2019structural,tenney2019bert,jawahar2019does}, and
attention-head analyses identify heads specialised for structural
functions~\citep{voita2019analyzing}. Attention sinks, where attention
mass concentrates on a few short structural tokens, were named by
\citet{xiao2024efficient} and independently observed in vision
Transformers~\citep{darcet2024registers}; later work links sink emergence
to residual-stream spectral filtering and training-time optimiser
dynamics~\citep{cancedda2024spectral,gu2025attention}. Our geometric
account explains this token preference quantitatively: function tokens
project onto the leading singular directions of $\Wdown$ differently from
content words, matching the Top-$1$ MA concentration observed in RQ3.

\paragraph{SVD geometry and causal intervention.}
SVD analyses of Transformer weights have mostly served compression, such
as activation-aware low-rank decomposition for models with outlier
activations~\citep{yuan2023asvd}. Causal-intervention methods such as
ROME~\citep{meng2022locating}, grounded in Pearl's framework for
causality~\citep{pearl2009causality}, support ablation-based attribution,
and automated circuit discovery extends such interventions to larger
subgraphs~\citep{conmy2023towards}. We combine these strands with a new
$V$-matrix ablation and a closed-form prediction (Eq.~\eqref{eq:delta-v}),
providing a quantitative bridge between SVD geometry and the generation
pathway of pathological activations.

% ============================================================
% 3. Preliminaries
% ============================================================
\section{Problem Formulation}
\label{sec:prelim}

Let $\mathbf{A} \in \mathbb{R}^{B \times L \times D}$ denote an internal
activation tensor with batch size $B$, sequence length $L$, and hidden
dimension $D$. For the MLP at layer $\ell$ we write the down-projection as
$\Wdown \in \mathbb{R}^{D \times \ffd}$ with singular decomposition
$\Wdown = \mathbf{U}\boldsymbol{\Sigma}\mathbf{V}^{\top}$, singular values
$\sigma_1 \geq \sigma_2 \geq \cdots \geq 0$, and columns
$\mathbf{u}_i \in \mathbb{R}^D$, $\mathbf{v}_i \in \mathbb{R}^{\ffd}$. The
MLP intermediate representation that feeds $\Wdown$ is denoted
$\mathbf{h}_2 \in \mathbb{R}^{\ffd}$.

\subsection{Definitions}
\label{sec:defs}

\paragraph{Definition 1 (Massive activation).}
A value $a \in \mathbf{A}$ is a \emph{massive activation} (MA) if $|a| > T$
where $T = P_{0.999}(|\mathbf{A}|)$, the $99.9$th percentile of absolute
activations in the tensor. We write
$\mathcal{M} = \{a \in \mathbf{A} : |a| > T\}$. This percentile rule is an
operational candidate definition; claims about orders-of-magnitude scale refer
to the measured Top-$1$ values rather than following from the percentile alone.

\paragraph{Definition 2 (Top-1 activation).}
$\topk_1 = \max_{b,l,d} |A_{b,l,d}|$.

\paragraph{Definition 3 ($\Delta\topk_1$ change rate).}
For an intervention $\mathcal{I}$,
$\Delta\topk_1 = (\topk_1^{(\mathcal{I})} - \topk_1^{(\text{base})}) / \topk_1^{(\text{base})} \times 100\%$.

\paragraph{Definition 4 (SVD alignment).}
For a vector $\mathbf{a} \in \mathbb{R}^{\ffd}$ and a singular direction
$\mathbf{v}_i$,
\(
  \varrho(\mathbf{a}, \mathbf{v}_i) = \mathbf{a}^{\top} \mathbf{v}_i / (\|\mathbf{a}\|_2 \|\mathbf{v}_i\|_2) \in [-1,1].
\)
Values $|\varrho| \to 1$ indicate strong directional alignment.

\paragraph{Definition 5 (Spectral dominance).}
$\eta = \sigma_1 / \sigma_2$. Larger $\eta$ means a more spectrally
concentrated $\Wdown$.

\paragraph{Definition 6 (Function token).}
A token $t$ is a \emph{function token} (FT) if, after stripping a leading
byte-pair-encoding space marker, its normalised form belongs to the union
$\mathcal{F} = \mathcal{F}_{\text{paper}} \cup \mathcal{F}_{\text{struct}}
\cup \mathcal{F}_{\text{digit}} \cup \mathcal{F}_{\text{bpe-frag}}$.
These four sets respectively cover closed-class function words, structural
tokens, digit pieces, and short BPE fragments; the complete word list,
regular-expression patterns, and detection procedure are given in
Appendix~\ref{app:ft-definition}.

\paragraph{Definition 7 (MA regime).}
Each model is assigned to one of four regimes based on the cumulative
layer-wise distribution of its MA intensity (\S\ref{sec:rq2}):
\textsc{Concentrated} if one MLP layer carries $\geq\!80\%$,
\textsc{Few-Source} if $2$--$4$ layers are required,
\textsc{Dispersed} if $\geq\!5$ layers are required, and
\textsc{Anomaly} if MLP ablation leaves $>\!50\%$ of MA intact. Spectral
dominance (Def.~5) and alignment (Def.~4) are measured consequences of
these regimes rather than defining criteria; operational details appear in
Appendix~\ref{app:models}.

\subsection{Research questions}
\label{sec:rq-list}

We structure our analysis around six linked roles, each evaluated in
\S\ref{sec:exp}: \textbf{RQ1 Source}, attention regulates rather than
generates MAs; \textbf{RQ2 Localization}, the MLP down-projection supplies
the physical substrate; \textbf{RQ3 Trigger}, function tokens preferentially
activate the mechanism; \textbf{RQ4 Mechanism}, the SVD expansion formula
predicts MA magnitude; \textbf{RQ5 Causality}, $V$-matrix geometry is
load-bearing; and \textbf{RQ6 Sufficiency}, single-layer recovery tests
residual-stream dependence.

The full experimental setup (model inventory, data sampling protocol,
statistical conventions) is reported in Appendix~\ref{app:setup}.

% ============================================================
% 4. Method
% ============================================================
\section{Methodology}
\label{sec:method}

For each model we identify one or more \emph{trigger layers} at which
$\topk_1$ exceeds the median $\topk_1$ of its neighbouring layers by at least
$30\times$; at boundary layers we use the available neighbouring layer only.
At such a layer,
$\mathbf{y} = \Wdown \mathbf{h}_2 + \mathbf{b}_{\text{down}} \in \mathbb{R}^{D}$
is the MLP output just before the residual stream and
$j^{*} = \arg\max_j |y_j|$ is the hidden dimension that carries the MA. The
down-projection bias $\mathbf{b}_{\text{down}}$ is identically zero on
LLaMA/Qwen-style architectures but nonzero on GPT-J, BLOOM, OPT, and GPT-2.

\paragraph{Trigger-layer terminology.}
We denote the MA manifestation layer by $L_{\text{surge}}$, distinct from
the first precursor layer $L_{\text{seed}}$, the maximum-ablation layer
$L_{\text{ampl}}$, and the peak-magnitude layer $L_{\text{peak}}$. RQ4 fits
Eq.~\eqref{eq:ma-poly} at the listed $L_{\text{surge}}$, while RQ5 ablates
the MLP down-projection at the same listed trigger layer, i.e. the block whose
output writes the extreme coordinate used in Table~\ref{tab:master}.

\subsection{SVD expansion formula}
\label{sec:method-formula}

By the SVD $\Wdown = \sum_{i=1}^{r} \sigma_i \mathbf{u}_i \mathbf{v}_i^{\top}$
with $r = \text{rank}(\Wdown)$, each coordinate of the MLP output admits the
exact expansion
\begin{equation}
  \label{eq:exact}
  y_j \;=\; (\Wdown \mathbf{h}_2 + \mathbf{b}_{\text{down}})_j \;=\; \sum_{i=1}^{r} \sigma_i\,(\mathbf{v}_i^{\top}\mathbf{h}_2)\,u_i[j] \;+\; (\mathbf{b}_{\text{down}})_j.
\end{equation}
Restricting to the MA channel $j^{*}$ gives our central generation formula
\begin{equation}
  \label{eq:ma-poly}
  \ma_{j^{*}} \;=\; \sum_{i=1}^{r} \sigma_i\,(\mathbf{v}_i^{\top}\mathbf{h}_2)\,u_i[j^{*}] \;+\; (\mathbf{b}_{\text{down}})_{j^{*}}.
\end{equation}
Eq.~\eqref{eq:ma-poly} is an \emph{algebraic identity} when summed to
$K=r$ (no approximation; it holds for any $\Wdown$). The falsifiable
empirical claim, tested in \S\ref{sec:rq4}, is that a low-rank
truncation at $K\leq 20$ recovers most of the MA magnitude. On modern
bias-free architectures $(\mathbf{b}_{\text{down}})_{j^{*}} = 0$ and the
formula collapses to a pure SVD expansion, while on legacy architectures the
bias term contributes a small additive offset (empirically $\sim 10\%$ of
the MA magnitude on GPT-J; \S\ref{sec:rq5}).

\paragraph{Three concurrent conditions for an MA.}
Each summand of Eq.~\eqref{eq:ma-poly} is a product of three factors, all
of which must be large for an MA to emerge at channel $j^{*}$:
(i) \emph{spectral power} $\sigma_i$ on the $i$-th singular direction;
(ii) \emph{directional alignment} $|\mathbf{v}_i^{\top}\mathbf{h}_2|$
between the MLP intermediate and that direction; and (iii) \emph{output
sparsity} $|u_i[j^{*}]|$, the concentration of the left-singular vector on
the MA channel. Failure of any single factor suppresses that term, and the
function-token preference of \S\ref{sec:rq3} traces directly to
factor (ii) at FT positions.

The empirical question is how few leading singular terms are needed to
recover the observed MA. We truncate at order $k$,
\begin{equation}
  \label{eq:ma-trunc}
  \widehat{\ma}^{(k)}_{j^{*}} \;=\; \sum_{i=1}^{k} \sigma_i\,(\mathbf{v}_i^{\top}\mathbf{h}_2)\,u_i[j^{*}] \;+\; (\mathbf{b}_{\text{down}})_{j^{*}},
\end{equation}
and measure the relative truncation error
$e^{(k)} = |\ma_{j^{*}} - \widehat{\ma}^{(k)}_{j^{*}}| / |\ma_{j^{*}}|$. The
sequence $k \mapsto e^{(k)}$ characterises how ``concentrated'' MA generation
is along the leading singular directions.

Flat-spectrum models are handled by increasing $k$: when no single
direction dominates, multiple leading terms can add coherently rather than
cancel. For \textsc{Few-Source} and \textsc{Dispersed} regimes we also use a
macro-SVD variant over the cumulative cross-layer residual delta. The
direction-consistency condition, $k$-selection rule, and macro-SVD formula
are given in Appendix~\ref{app:formula-extensions}; the empirical
$k \mapsto e^{(k)}$ curve is reported in Table~\ref{tab:polyfit}.

\subsection{Single-direction regression form}
\label{sec:method-reg}

When a single singular direction dominates---i.e.\ $\sigma_1 \gg \sigma_2$
and $|\varrho(\mathbf{h}_2, \mathbf{v}_1)| \to 1$, the spectrally dominated
sub-case of the \textsc{Concentrated} regime---Eq.~\eqref{eq:ma-trunc} at
$k{=}1$ reduces to a linear relationship between the observed MA and the
scalar projection $p(x) = \mathbf{v}_1^{\top}\mathbf{h}_2(x)$:
\begin{equation}
  \label{eq:lin-reg}
  \ma(x) \;\approx\; \beta_{\text{slope}}\,p(x) + b, \qquad \beta_{\text{slope}} = \sigma_1\,u_1[j^{*}],
\end{equation}
with an intercept
$b = \sum_{i \geq 2} \sigma_i\,(\mathbf{v}_i^{\top}\mathbf{h}_2)\,u_i[j^{*}] + (\mathbf{b}_{\text{down}})_{j^{*}} + r_{\text{stream}}$
that absorbs higher-order singular contributions, the down-projection bias,
and the residual-stream carry-over $r_{\text{stream}}$. The exact identity
of Eq.~\eqref{eq:ma-poly} applies to the MLP contribution at the trigger
layer; comparisons against residual-stream Top-$1$ activations therefore
use Eq.~\eqref{eq:lin-reg}, whose intercept absorbs the residual
carry-over. Because $\beta_{\text{slope}}$ and $b$ are jointly identifiable
from token-level observations, Eq.~\eqref{eq:lin-reg} can be fit directly
and evaluated by $R^2$ without separately measuring $\sigma_1$ or
$u_1[j^{*}]$.

\subsection{Closed-form causal prediction under random $\mathbf{V}$}
\label{sec:method-vcf}

Replace the right-singular matrix $\mathbf{V}$ by a Haar-uniform orthogonal
$\tilde{\mathbf{V}} \in \mathbb{R}^{\ffd \times \ffd}$ while keeping
$\mathbf{U}$, $\boldsymbol{\Sigma}$, and all other parameters (including
the normalisation affine parameters) fixed. In Eq.~\eqref{eq:lin-reg} the
only factor that changes is the projection, now
$\tilde{p} = \tilde{\mathbf{v}}_1^{\top}\mathbf{h}_2$.

By concentration on the sphere, $\tilde{p}$ is approximately half-normal in
magnitude with mean $\sqrt{2/\pi}\|\mathbf{h}_2\|_2/\sqrt{\ffd}$. Combining
this with $|p| \approx \|\mathbf{h}_2\|_2$ in the spectrally dominated
sub-case of the \textsc{Concentrated} regime gives the expected signed change
\begin{equation}
  \label{eq:delta-v}
  \mathbb{E}[\Delta_V] \;\approx\; \frac{\sqrt{2/\pi}}{\sqrt{\ffd}} - 1.
\end{equation}
For Qwen2.5-7B ($\ffd = 18{,}944$), this predicts $-99.42\%$, matching the
empirical $-99.2\%$ in RQ5 to within $0.25$ percentage points. A formal
derivation with finite-$\ffd$ corrections is given in
Appendix~\ref{app:proofs}.

\subsection{Intervention algorithms}
\label{sec:method-interventions}

We use four $V$-matrix interventions: Haar-random replacement of
$\mathbf{V}$, projection removal of the leading $K$ right-singular
directions, macro-$V$ ablation for multi-layer regimes, and per-expert
ablation for MoE models. A model passes RQ5 if any D1--D4 criterion is met
(single-layer multi-$K$, per-dimension collapse, macro collapse, or a
two-point boundary relaxation). Per-model ablation magnitudes appear in
Table~\ref{tab:master}; full algorithms and thresholds are in
Appendix~\ref{app:interventions}.

\paragraph{Regime classification.}
Regime assignment follows Def.~7. In \S\ref{sec:rq4}--\S\ref{sec:rq5} we
show that these assignments are consistent across the full 26-model pool,
with the signed $V$-ablation change tracking the truncation error $e^{(k)}$ as
predicted by Eq.~\eqref{eq:delta-v}.

% ============================================================
% 5. Experiments
% ============================================================
\section{Experiments and Results}
\label{sec:exp}

Each research question follows the schema \emph{Hypothesis} $\rightarrow$
\emph{Intervention} $\rightarrow$ \emph{Result} $\rightarrow$
\emph{Interpretation}. Table~\ref{tab:master} collects the per-model
summary used throughout the experiments; the remainder of this section
surfaces regime-level findings that match the theoretical predictions of
\S\ref{sec:method}.

\begin{table}[t]
  \caption{Complete 26-model dashboard. The table reports, for every model,
  the MA regime, trigger layer, SVD-geometry indicators, causal $V$-ablation
  effects, and RQ1 attention mode. The dashboard shows only single- and
  macro-$V$ signed changes; the per-coordinate (D2) and boundary (D4)
  diagnostics that contribute to the full RQ5 D1--D4 verdict are not
  table-visible. N/A denotes interventions that were not run or not
  applicable; macro values are reported only when used by a diagnostic or
  exception criterion.}
  \label{tab:master}
  \centering
  \begingroup
  \scriptsize
  \setlength{\tabcolsep}{2.4pt}
  \renewcommand{\arraystretch}{0.82}
  \begin{tabular*}{\textwidth}{@{\extracolsep{\fill}}clccccccc@{}}
    \toprule
    \# & Model & Type & Trigger $L$ & Spec.\ $\eta$ & Fit $R^2$ &
    Single-$V$ $\Delta$ & Macro-$V$ $\Delta$ & RQ1 mode \\
    \midrule
    \multicolumn{9}{l}{\textit{\textsc{Concentrated} (single-layer MA origin; 8 models + 1 ambiguous)}} \\
    \midrule
    $\star$1  & GPT-J-6B             & Dense & 2  & 2.52 & 1.00 & $-99.0$ & $-99.1$ & Gen* \\
    $\star$2  & Qwen2-7B             & Dense & 3  & 2.84 & 1.00 & $-98.6$ & N/A     & Gen  \\
    $\star$3  & Qwen2.5-7B           & Dense & 3  & 2.64 & 1.00 & $-99.2$ & N/A     & Sup  \\
    $\star$4  & Qwen3-0.6B           & Dense & 2  & 1.41 & 1.00 & $-93.0$ & N/A     & Gen  \\
    5         & GLM4-32B             & Dense & 0  & 1.53 & 0.47 & $-96.8$ & $-17.1$ & Sup  \\
    6         & Mistral-7B-v0.3      & Dense & 1  & 1.12 & 1.00 & $-83.0$ & N/A     & Gen  \\
    7         & Qwen2.5-0.5B         & Dense & 2  & 1.48 & 0.91 & $-55.0$ & N/A     & Gen  \\
    8         & LLaMA-2-13B          & Dense & 0  & 1.32 & 0.97 & $-95.6$ & N/A     & Gen  \\
    9$^{\dagger}$ & LLaMA-2-7B-Chat  & Dense & 1  & 1.04 & 0.94 & $-95.7$ & N/A     & Sup  \\
    \midrule
    \multicolumn{9}{l}{\textit{\textsc{Few-Source} (2--4 trigger layers; 7 models)}} \\
    \midrule
    10  & GPT-2                      & Dense & 3  & 3.05 & 0.55 & $-6.5$  & $-94.6$  & Gen  \\
    11  & Qwen3-1.7B                 & Dense & 2  & 1.33 & 0.94 & $-86.8$ & $-99.8$  & Gen  \\
    12  & Qwen3-4B                   & Dense & 6  & 1.24 & 1.00 & $-94.9$ & $-99.7$  & Gen  \\
    13  & Falcon-7B                  & Dense & 3  & 1.37 & 0.99 & $-98.1$ & $-97.4$  & Gen  \\
    14  & LLaMA-3.1-8B               & Dense & 1  & 1.38 & 1.00 & $-93.0$ & $-99.8$  & Gen  \\
    15  & GLM4-9B                    & Dense & 1  & 3.26 & 0.89 & $-46.5$ & $-81.8$  & Gen  \\
    16  & BLOOM-7B1                  & Dense & 7  & 1.81 & 1.00 & $-69.7$ & $-82.0$  & Gen* \\
    \midrule
    \multicolumn{9}{l}{\textit{\textsc{Dispersed} ($\geq 5$ trigger layers; 8 models)}} \\
    \midrule
    17  & Qwen3-8B                   & Dense  & 6  & 1.48 & 1.00 & $-95.6$ & $-99.6$  & Gen  \\
    18  & Yi-9B                      & Dense  & 8  & 1.43 & 0.88 & $-93.1$ & $-99.2$  & Sup  \\
    19  & Qwen3.5-9B$^{\ddagger}$    & Hybrid & 22 & 1.06 & 0.73 & $-70.2$ & $-0.2$   & Gen  \\
    20$^{\ddagger\ddagger}$ & Qwen1.5-14B & Dense  & 2  & 1.33 & 1.00 & $-49.3$ & $-12.6$  & Gen  \\
    21  & Qwen3-14B                  & Dense  & 6  & 1.33 & 1.00 & $-93.5$ & $-88.2$  & Sup  \\
    22  & Qwen3.5-27B$^{\ddagger}$   & Hybrid & 54 & 1.12 & 0.99 & $-77.9$ & $-0.4$   & Gen  \\
    23  & Qwen3-30B-A3B$^{\ast}$     & MoE    & 1  & 1.17 & 0.38 & $-0.8$  & $-0.4$   & Gen  \\
    24  & Qwen3-32B                  & Dense  & 6  & 1.35 & 1.00 & $-97.9$ & $-86.3$  & Sup  \\
    \midrule
    \multicolumn{9}{l}{\textit{\textsc{Anomaly} (MLP retain $> 50\%$; 2 models)}} \\
    \midrule
    25  & OPT-6.7B                   & Dense  & 1  & 2.53 & 0.98 & $-17.5$ & N/A      & Sup  \\
    26  & Qwen3.5-35B-A3B$^{\ast,\ddagger}$ & MoE+H & 9  & 1.03 & 0.00 & $+0.1$  & $+0.5$   & Sup  \\
    \bottomrule
  \end{tabular*}
  \endgroup

  \smallskip
  \footnotesize
  \begin{flushleft}
    Models are grouped by MA regime (Def.~7): \textsc{Concentrated},
    \textsc{Few-Source}, \textsc{Dispersed}, and \textsc{Anomaly}.
    Type abbreviations: Dense, Hybrid attention, MoE, and MoE+H.
    $\star$ marks hero models that pass every RQ1--RQ5 test, where
    \emph{passing RQ1} means that attention is not the direct writer of
    the extreme coordinate (i.e.\ the MLP retains the trigger-layer source)
    rather than that attention removal leaves the MA intact.
    $\dagger$ marks the ambiguous RLHF layer-support case grouped with
    \textsc{Concentrated}. $\ddagger$ marks hybrid-attention bypasses;
    $\ast$ marks MoE routing dilution. $\ddagger\ddagger$ marks
    \emph{D2 per-coordinate diagnostic rescue}: although the mean
    Single-/Macro-$V$ entries do not cross the $-80\%$ threshold, the
    per-coordinate $\Delta\mathrm{MA}_{j^{*}} = -100\%$ at the MA channel
    yields a full RQ5 PASS under the D2 criterion (\S\ref{sec:rq5}).

    Macro-$V$ is reported as N/A for \textsc{Concentrated} models per
    the pre-registered category-bound criterion: \textsc{Concentrated}
    checkpoints are evaluated on D1 single-layer $V$-ablation, while
    macro-$V$ is consulted only when D1 falls in a boundary band or as
    an exception diagnostic.

    The dense pool of $22$ models used for the headline RQ1--RQ5 PASS
    rates excludes four architecture-anomaly checkpoints: OPT-6.7B
    (Tier~E, pre-LayerNorm + non-standard FFN), Qwen3.5-9B (Tier~C,
    hybrid-attention bypass), Qwen3.5-35B-A3B (Tier~C, MoE + hybrid),
    and Qwen3-30B-A3B (Tier~C, MoE routing dilution). These four are
    reported in this dashboard but excluded from the dense-pool
    aggregate ratios in \S\ref{sec:rq1}--\S\ref{sec:rq5}.
  \end{flushleft}
\end{table}

% ------------------------------------------------------------
\subsection{(RQ1) Source: attention regulates rather than generates MAs}
\label{sec:rq1}

\paragraph{Hypothesis and intervention.}
If attention \emph{generates} MAs, zeroing every attention output should
eliminate them. We zero all attention module outputs at every layer via
PyTorch forward hooks, leaving the residual stream and MLP untouched, and
measure $\Delta\topk_1$ at the baseline trigger layer.

\paragraph{Zeroing attention does not erase MAs.}
Across the twenty-six models, the residual ratio
$\rho = \topk_1^{(\text{no-attn})}/\topk_1^{(\text{base})}$ separates the
outcomes into three qualitatively distinct behaviours that cut across
model scale and family. In the first, which we label
\textbf{Gen}\textsuperscript{*}, attention carries essentially all of
the downstream amplification: removing it collapses the MA by more than
ninety-five percent, and only BLOOM-7B1 and GPT-J-6B exhibit this
strong dependence. The second and most common behaviour, \textbf{Gen},
covers sixteen models in which attention amplifies an MA that the MLP
has already begun to produce, so that attention removal reduces but
does not eliminate the signal. The third behaviour, \textbf{Sup}, covers
eight models in which attention acts as a stable suppressor: zeroing it
actually grows the MA, in several cases by more than a factor of two.
The resulting RQ1 modes and regime assignments are reported in
Table~\ref{tab:master}; full residual ratios are left to the supplement.

\paragraph{Interpretation.}
Because $\Delta\topk_1$ ranges from strongly negative to strongly
positive across the pool, no single ``attention-as-generator'' account
fits these data. Since substantial MA magnitude persists under many of
the interventions, attention is best understood as a \emph{regulator}
that gates or suppresses the downstream generator rather than as the
primary generator of the extreme coordinate. Even in Gen\textsuperscript{*}
cases, attention supplies downstream amplification while the coordinate is
written through the MLP pathway.

\paragraph{Key finding~1.}
\emph{Attention mechanisms act as regulators---in one of three observed
modes (Gen\textsuperscript{*}, Gen, Sup)---rather than as primary
generators of the MA substrate
of massive activations.}

% ------------------------------------------------------------
\subsection{(RQ2) Localization: MLP down-projection supplies the physical substrate}
\label{sec:rq2}

\paragraph{Hypothesis and intervention.}
At every trigger layer we (a) compare the attention-block output
$\mathbf{H}_{\text{attn}}$ against the MLP-block output
$\mathbf{H}_{\text{mlp}}$ via the dominance ratio
$\rho_\ell = \max|\mathbf{H}_{\text{mlp}}| / \max|\mathbf{H}_{\text{attn}}|$;
and (b) ablate the MLP at that layer, recording the retention ratio
\begin{equation}
  \label{eq:tau-retain}
  \tau \;=\; \frac{\topk_1^{(\text{mlp-zero})}}{\topk_1^{(\text{base})}}.
\end{equation}
A model passes RQ2 strictly when $\tau \leq 0.10$ and at the boundary when
$\tau \leq 0.15$. In parallel, $\rho_\ell > 1$ is inconsistent with equal
MLP and attention magnitudes in the observed measurement at that layer.

\paragraph{Removing the MLP removes the spike.}
Across dense models the MLP output exceeds the attention output by
between a factor of a few and more than three orders of magnitude,
confirming that the MLP branch dominates the residual stream at every
trigger layer. Under single-layer MLP ablation, twenty-one of the
twenty-four non-MoE checkpoints retain less than ten percent of their
baseline Top-$1$ activation, directly implicating $\Wdown$ as the
primary substrate of the MA. The remaining three non-MoE models fall
outside this ``MLP-dominant'' pattern in two different ways. OPT-6.7B
(together with the MoE variant Qwen3.5-35B-A3B, whose expert router
dilutes the MA signal) crosses the fifty-percent retain threshold and
is therefore classified as \textsc{Anomaly} in Def.~7. Qwen3.5-9B,
whose hybrid linear-attention branch partially bypasses $\Wdown$, and
GLM4-32B occupy an intermediate regime in which $\Wdown$ dominates but
not exclusively; on both, the single-layer $V$-ablation still drives
the MA down by more than seventy percent. We return to these
architectural causes in \S\ref{sec:moe-dilution}; the aggregate RQ2
outcome is summarised in Table~\ref{tab:polyfit}.

\paragraph{Key finding~2.}
\emph{Across $23/26$ models ($88.5\%$ under the boundary criterion
$\tau \leq 0.15$), the MLP---specifically the down-projection
$\Wdown$---is responsible for most MA intensity, with the failures
concentrated in routing, bypass, or anomaly cases.}

% ------------------------------------------------------------
\subsection{(RQ3) Trigger: function tokens preferentially activate the mechanism}
\label{sec:rq3}

\paragraph{Hypothesis and intervention.}
We assign FT labels using Def.~6, which combines a hand-curated
function-token list with structural-token, digit-token, and BPE-fragment
rules. For each model we record the token associated with
$\topk_1$ at the trigger layer across the $64$ C4 windows, and compute,
for $K \in \{1, 5, 10, 20\}$, the FT trigger rate
\begin{equation}
  \label{eq:pi-func}
  \pi_{\text{func}} \;=\; \frac{n_{11}}{n_{+1}},
\end{equation}
where $n_{11}$ counts trigger-layer Top-$K$ MA tokens that are FTs and
$n_{+1}$ is the total number of Top-$K$ measurements. To test whether the
FT enrichment is statistically distinguishable from independence we run
Fisher's exact test on the $2\!\times\!2$ contingency table of (token is
FT) $\times$ (token triggers MA). The probability of the observed table is
\begin{equation}
  \label{eq:fisher-table}
  P(T=n_{11}) \;=\; \frac{\binom{n_{1+}}{n_{11}}\binom{n_{2+}}{n_{21}}}{\binom{N}{n_{+1}}},
\end{equation}
and the Fisher $p$-value sums this probability over tables at least as
extreme, declaring significance at $p < 0.01$.

\paragraph{Function tokens are the trigger points.}
The Top-$1$ MA token is a function token in twenty-four of the
twenty-six models, well above the same-corpus baseline FT rate of
roughly one in three. Enrichment at Top-$5$ and Top-$10$ remains
substantially above baseline and decays toward it as the rank $K$
grows, consistent with a mechanism sharply localised on the very
highest-magnitude positions rather than on the broader tail. On the
\textsc{Concentrated} models for which a second-direction diagnosis is
informative---excluding the rank-$1$-exact GPT-J-6B and the ambiguous
LLaMA-2-7B-Chat case---the leading two singular directions
$\mathbf{v}_1$ and $\mathbf{v}_2$ place their strongest projection on
the same function token in all but one case, with the shared token
typically a structural whitespace or punctuation marker (examples
include \texttt{\textbackslash n\textbackslash n}, \texttt{`~@'}, and
\texttt{`S'}). This direction coherence explains why the SVD-expansion
rescue of RQ4 preserves the FT interpretation as the truncation order
$k$ increases.

\paragraph{Decoding $\mathbf{u}_1$ through the unembedding.}
The left singular vector $\mathbf{u}_1$ tells us \emph{which output
vocabulary token} the MA channel writes into. Multiplying it through the
unembedding matrix $\mathbf{W}_U$ gives a logit vector over the vocabulary
$\mathcal{V}$:
\begin{equation}
  \label{eq:u1-decode}
  \text{logits}_{\mathbf{u}_1} \;=\; \mathbf{W}_U\,\mathbf{u}_1,
  \qquad \mathrm{Top}\text{-}1_{\mathbf{u}_1} \;=\; \arg\max_{t \in \mathcal{V}} \bigl(\text{logits}_{\mathbf{u}_1}[t]\bigr).
\end{equation}
On six of these seven \textsc{Concentrated} models, the
$\mathrm{Top}\text{-}1_{\mathbf{u}_1}$ token is the \emph{same} structural
FT (typically \texttt{`\textbackslash n\textbackslash n'} or another
whitespace/punctuation marker) that achieves the largest input-side
projection on $\mathbf{v}_1$. The MA channel therefore closes a loop:
the FT both triggers ($\mathbf{v}_1$) and is amplified ($\mathbf{u}_1$) by
the same MLP geometry.

\paragraph{Frequency-vs-syntax decoupling.}
\label{sec:freq-vs-syntax}
Because FTs are inherently high-frequency, one might suspect the
enrichment is a pure frequency effect. We disentangle category from
frequency along three converging lines.

\emph{(i) $2\!\times\!2$ contingency at the trigger layer.}
Define the high-frequency set
$\mathcal{H}_{\theta} = \{\,t \in \mathcal{V}: \mathrm{freq}(t) \geq \theta\,\}$
with $\theta$ chosen so that $|\mathcal{H}_{\theta}| = K$. Crossing this
with the function-token set $\mathcal{F}$ yields four quadrants
$Q_1{=}\mathcal{F}\cap\mathcal{H}$,
$Q_2{=}\mathcal{F}\setminus\mathcal{H}$,
$Q_3{=}\mathcal{H}\setminus\mathcal{F}$,
$Q_4{=}\overline{\mathcal{F}}\cap\overline{\mathcal{H}}$, and a
within-quadrant trigger rate
$\pi(Q_i) = |\{\,t\in Q_i : t \text{ triggers MA}\,\}| / |Q_i|$.
The frequency-confound hypothesis predicts
$\pi(Q_2) \!\approx\! \pi(Q_4)$ and $\pi(Q_3) \!\approx\! \pi(Q_1)$;
the syntax hypothesis predicts the opposite asymmetry,
$\pi(Q_2)\!\gg\!\pi(Q_3)$. Table~\ref{tab:freq-vs-syntax} reports the
GPT-2 measurement at trigger layer $L{=}3$. Low-frequency function
tokens trigger MA at rate $0.612$ while high-frequency content tokens
trigger at only $0.014$:
\begin{equation}
  \label{eq:freq-vs-syntax-ratio}
  \frac{\pi(Q_2)}{\pi(Q_3)} \;=\; \frac{0.612}{0.014} \;\approx\; 43.7,
\end{equation}
a $43.7\!\times$ separation in favour of the syntactic category.

\emph{(ii) Logistic regression with frequency as a covariate.}
We fit
\begin{equation}
  \label{eq:logit-freq-syntax}
  \mathrm{logit}\bigl(P(t \text{ triggers MA})\bigr)
  \;=\; \beta_0 + \beta_{\mathcal{F}}\,\mathbb{1}[t\in\mathcal{F}]
       + \beta_{\text{freq}}\,\log\mathrm{freq}(t),
\end{equation}
with document-cluster-robust ($G{=}30$) sandwich SE. Across the $22$
dense checkpoints we obtain $\hat{\beta}_{\mathcal{F}} > 0$ at
$z > 3.29$ ($p < 10^{-3}$) on every model and
$|\hat{\beta}_{\mathcal{F}}| > 2|\hat{\beta}_{\text{freq}}|$ in $20/22$,
satisfying our pre-registered criterion (Appendix~\ref{app:ft-definition}).
On Qwen2.5-7B, $\hat{\beta}_{\mathcal{F}} = 1.83$ ($p < 10^{-12}$).

\emph{(iii) Permutation null and additional support.}
Shuffling FT labels uniformly at random ($B{=}1000$) yields a null
distribution for the ratio $\pi(Q_2)/\pi(Q_3)$ with mean $1.02$ and
$95\%$ percentile interval $[0.83, 1.21]$; the observed value $43.7$
lies far outside this band ($p < 10^{-3}$). Two further observations
reinforce the same conclusion: the causal $V$-ablation of RQ5 leaves
all input-side statistics unchanged yet still suppresses $78$--$99\%$
of MA on D1--D4-PASS models, and MAs are layer-specific (the same FT
does not produce an MA at non-trigger layers).

\begin{table}[t]
  \caption{Frequency-vs-syntax $2\!\times\!2$ on GPT-2 at trigger layer
  $L{=}3$ (WikiText-2, $30$ documents, $K{=}1000$). Trigger threshold:
  $99$th-percentile MA over the window. Low-frequency FTs trigger MA
  $43.7\!\times$ more often than high-frequency content tokens.}
  \label{tab:freq-vs-syntax}
  \centering
  \small
  \begin{tabular}{@{}l c c c@{}}
    \toprule
    Quadrant & Example tokens & $|Q_i|$ & $\pi(Q_i)$ \\
    \midrule
    $Q_1=\mathcal{F}\cap\mathcal{H}$ (HF FT)
      & \texttt{the}, \texttt{of}, \texttt{.} & $156$ & $0.795$ \\
    $Q_2=\mathcal{F}\setminus\mathcal{H}$ (LF FT)
      & \texttt{`\textbackslash n\textbackslash n'}, \texttt{`\,k'}, \texttt{St} & $1{,}243$ & $\mathbf{0.612}$ \\
    $Q_3=\mathcal{H}\setminus\mathcal{F}$ (HF content)
      & \texttt{language}, \texttt{model}, \texttt{data} & $287$ & $\mathbf{0.014}$ \\
    $Q_4=\overline{\mathcal{F}}\cap\overline{\mathcal{H}}$ (LF content)
      & rare proper nouns & $8{,}952$ & $0.011$ \\
    \bottomrule
  \end{tabular}
\end{table}

\paragraph{Function tokens are information-theoretic anchors.}
Beyond raw frequency, FT positions are distinguished by the \emph{low
entropy} of their next-token prediction. On $14$ models sampled for
detailed conditional-entropy analysis, FT positions exhibit a
next-token entropy
$H(C) = -\sum_{c} p(c \mid x) \log p(c \mid x)$ consistently $2$--$4$
bits below content-token positions (e.g.\ on Qwen2.5-7B a
\texttt{`\textbackslash n\textbackslash n'} position gives
$H(C) \approx 0.3$ bits, whereas a typical content-word position gives
$H(C) \approx 3.5$ bits). This \emph{information-theoretic anchor}
interpretation explains why MA concentrates on FTs rather than on
arbitrary high-frequency tokens: the language model has the tightest
prediction context at low-entropy positions, and our geometric
mechanism shows those positions are precisely the ones whose
$\mathbf{h}_2$ aligns with $\mathbf{v}_1$. Function tokens are thus
simultaneously sparse in two complementary senses---token-level
(low-entropy anchors) and direction-level ($\mathbf{v}_1$,
$\mathbf{u}_1$ spectrum)---a \emph{dual-sparsity} framing we develop
in \S\ref{sec:dual-sparsity}.

\paragraph{Key finding~3.}
\emph{Function tokens act as geometric anchors of MA generation: they
project onto the leading singular directions of $\Wdown$ in a way that
content words do not, and they do so coherently across the top two
singular directions in \textsc{Concentrated} models.}

% ------------------------------------------------------------
\subsection{(RQ4) Mechanism: the SVD expansion predicts MA magnitude}
\label{sec:rq4}

\paragraph{Hypothesis and intervention.}
Eq.~\eqref{eq:ma-poly} predicts $\ma_{j^{*}}$ as a truncated sum over the
leading $k$ singular directions. We compute the relative truncation error
$e^{(k)}$ for the trigger layer of each model over a $1{,}000$-token
WikiText-2 window, at $k \in \{1, 3, 5, 10, 20\}$.

\begin{table}[t]
  \caption{Compact evidence matrix for the six-question mechanism chain.
  The top block reports aggregate results; the bottom block keeps the RQ4
  rank-$k$ truncation curve in the main text.}
  \label{tab:polyfit}
  \centering
  \scriptsize
  \setlength{\tabcolsep}{4pt}
  \renewcommand{\arraystretch}{0.9}
  \begin{tabular}{@{}l l c l@{}}
    \toprule
    RQ & Main test & Headline statistic & Mechanistic role \\
    \midrule
    RQ1 & attention zeroing & Gen$^{*}$: $2$, Gen: $16$, Sup: $8$ & regulator \\
    RQ2 & MLP ablation & $23/26$ pass under $\tau \leq 0.15$ & $\Wdown$ substrate \\
    RQ3 & FT enrichment & $24/26$ Top-$1$ FT triggers & function-token anchors \\
    RQ4 & SVD expansion & $24/26$ unified; rank-$20$ $16/22$ & magnitude prediction \\
    RQ5 & D1--D4 $V$-ablation & $21/26$ full; dashboard threshold $18/21$ & causal geometry \\
    RQ6 & single-layer recovery & $2/26$ high; dense $16/22$ consistent & residual stream \\
    \bottomrule
  \end{tabular}

  \vspace{2pt}
  \begin{tabular}{@{}l c c c c c@{}}
    \toprule
    RQ4 rank $k$ & $1$ & $3$ & $5$ & $10$ & $20$ \\
    \midrule
    near-exact ($e^{(k)}<5\%$) & $6$ & $6$ & $6$ & $7$ & $\mathbf{9}$ \\
    PASS ($e^{(k)}<30\%$) & $11/22$ & $15/22$ & $15/22$ & $15/22$ & $\mathbf{16/22}$ \\
    PASS rate & $50\%$ & $68\%$ & $68\%$ & $68\%$ & $\mathbf{73\%}$ \\
    \bottomrule
  \end{tabular}
\end{table}

\paragraph{Low-rank SVD predicts the spike.}
Table~\ref{tab:polyfit} shows that a single singular direction already
predicts the observed MA within thirty percent on half of the
twenty-two models analysed, and within five percent on six of them.
Extending the truncation to $k{=}20$ raises the pass rate to roughly
seventy percent, with nine models in near-exact agreement. Four hero
\textsc{Concentrated} models---GPT-J-6B, Qwen2.5-7B, Qwen2-7B, and
Qwen3-0.6B---achieve $R^{2} = 1.00$ and essentially zero truncation
error at rank one, making them the cleanest empirical instantiation of
Eq.~\eqref{eq:ma-poly}. The models that benefit most from higher
truncation orders are, by inspection of Table~\ref{tab:master},
precisely those with flatter spectra ($\eta < 2$), so the SVD
truncation error tracks spectral structure in the way
Eq.~\eqref{eq:ma-poly} predicts.

\paragraph{Empirical verification of $\beta_{\text{slope}} = \sigma_1 \cdot u_1[j^{*}]$.}
The regression form of Eq.~\eqref{eq:lin-reg} identifies the fitted
slope $\beta_{\text{slope}}$ with the product $\sigma_1 \cdot u_1[j^{*}]$,
a mechanistic identity that we test by comparing the regression coefficient against
the SVD-derived factor. On the two hero Qwen models, Qwen2-7B and
Qwen2.5-7B, the fitted slope agrees with the predicted product within
two percent, directly confirming the identity. More broadly, the
leading coordinate of $\mathbf{u}_1$ captures at least half of its
total energy on thirteen of the fifteen models for which the full
$\Wdown$ SVD is tabulated, with one near-degenerate extreme (GPT-2)
and two dense counter-examples (Qwen3-8B and Qwen3-14B).
The $\mathbf{u}_1$ sparsity complements the input-side
$\mathbf{h}_2$--$\mathbf{v}_1$ alignment: together they localise MA to
a single hidden coordinate $j^{*}$ exactly as Eq.~\eqref{eq:lin-reg}
predicts.

\paragraph{Key finding~4.}
\emph{Low-rank truncation of Eq.~\eqref{eq:ma-poly} at $K\leq 20$
empirically captures the observed MA magnitude on most models:
the rank-$20$ single-layer truncation, evaluated at
$k \in \{1, 3, 5, 10, 20\}$, recovers the observed MA within $30\%$\footnote{The
$30\%$ threshold corresponds to errors below the random-$K$ null
$95$th percentile (Appendix~\ref{app:formula-extensions}). Top-$K$
truncation errors on the passing models are typically $5$--$20\times$
tighter than random-$K$ baselines.} on
$16/22$ models and within $5\%$ on $9/22$; the rank-$1$ truncation alone
is sufficient for half of the $22$ analysed models and is near-exact
for six. Combining single-layer multi-$K$ $V$-projection (the D1
criterion of \S\ref{sec:method-interventions} at
$K \in \{1, 3, 10, 20\}$) with the macro-SVD path, the unified SVD
criterion ($K \in \{1, 3, 10, 20\}$ or macro $\Delta_V \leq -80\%$) is
satisfied by $24/26 = 92.3\%$ of all models, with the two failures
attributable to routing dilution in MoE architectures and
hybrid-attention bypass.}

% ------------------------------------------------------------
\subsection{(RQ5) Causality: $V$-matrix geometry is load-bearing}
\label{sec:rq5}

\paragraph{Hypothesis and intervention.}
If Eq.~\eqref{eq:ma-poly} is causal rather than merely predictive, then
randomising the right-singular matrix while preserving
$\boldsymbol{\Sigma}$, $\mathbf{U}$, and all other parameters should collapse
MA magnitude by an amount governed by Eq.~\eqref{eq:delta-v}.
We perform the single-layer intervention of
\S\ref{sec:method-interventions} on every model and, for
non-\textsc{Concentrated} models, the macro variant.

\paragraph{Randomising $V$ collapses the mechanism.}
Twenty-one of the twenty-six models pass the combined D1--D4 criterion of
\S\ref{sec:method-interventions} ($80.8\%$). The dashboard
(Table~\ref{tab:master}) shows only single-layer and macro-$V$ signed
changes; under this dashboard-visible threshold ($\Delta_V \leq -80\%$),
the clean dense subset excluding MoE, hybrid-attention, and the OPT
anomaly passes in $18/21$ cases. The full D1--D4 verdict additionally
counts the per-coordinate D2 and boundary D4 diagnostics, which are not
table-visible and rescue three further models. Among
the passing \textsc{Concentrated} models the closed-form prediction of
Eq.~\eqref{eq:delta-v} is tight: on Qwen2.5-7B the empirical single-layer
collapse matches the theoretical prediction to within a quarter of a
percentage point. \textsc{Dispersed} models depart from this prediction in the
expected direction, with the gap tracking the RQ4 truncation error
$e^{(1)}$; \textsc{Few-Source} single-layer failures are rescued by macro
$V$-ablation, while MoE and hybrid-attention failures reflect routing
dilution or architectural bypasses rather than a dense $\Wdown$ mechanism.

\paragraph{Key finding~5.}
\emph{$V$-matrix geometry is load-bearing for MA generation in the
full pool (ablating it is sufficient for elimination): the combined D1--D4 criterion is satisfied by
$21/26 = 80.8\%$ of models, while the dashboard-visible single/macro-$V$
threshold ($\Delta_V \leq -80\%$) is met by $18/21$ clean dense models
after excluding MoE, hybrid-attention, and the OPT anomaly. The
magnitude of the collapse quantitatively matches the closed-form prediction
of Eq.~\eqref{eq:delta-v} on \textsc{Concentrated} models and deviates on
\textsc{Dispersed} models in proportion to the RQ4 truncation error
$e^{(1)}$.}

% ------------------------------------------------------------
\subsection{(RQ6) Sufficiency: single-layer recovery tests residual-stream dependence}
\label{sec:rq6}

\paragraph{Hypothesis and intervention.}
RQ2 establishes that the MLP block is \emph{necessary} for MA
generation. RQ6 tests the complementary sufficiency question: when
all other MLP outputs are zeroed along the forward pass, does
restoring a single MLP layer suffice to regenerate the MA? For each
layer $L$ we run three forward passes---\emph{baseline} (all MLPs
active), \emph{floor} (every MLP output zeroed), and
\emph{keep-$L$} (only layer $L$'s MLP restored)---and compute the
recovery rate
\begin{equation}
  \label{eq:recovery}
  r(L) \;=\; \frac{\topk_1^{(\text{keep-}L)} - \topk_1^{(\text{floor})}}
                  {\topk_1^{(\text{base})} - \topk_1^{(\text{floor})}} \times 100\%.
\end{equation}
The best single-layer recovery is $r^{*} = \max_{L} r(L)$, typically
achieved at the regime-specific trigger layer.

\paragraph{One restored MLP is rarely enough.}
Only $2/26$ models reach $r^{*} \geq 30\%$: GPT-J-6B at $76.4\%$
(trigger layer $L{=}2$) and LLaMA-3.1-8B at $49.0\%$ ($L{=}1$). Three
hero \textsc{Concentrated} models that passed RQ4 with $R^{2}=1.00$
nevertheless recover only $0$--$12\%$ under single-layer restoration
(Qwen2.5-7B: $-0.2\%$; Qwen2-7B: $0.3\%$; Qwen3-0.6B: $12.1\%$). The
collapse is \emph{not} a failure of the SVD expansion: the
formula's premise, that $\mathbf{h}_2$ aligns with $\mathbf{v}_1$, is
broken by the intervention itself. Zeroing every upstream MLP disrupts
the residual stream, so layer $L$'s MLP now sees a corrupted
$\mathbf{h}_2$ input and cannot regenerate the MA even though its
$\Wdown$ geometry remains intact.

\paragraph{Key finding~6.}
\emph{MA is a \emph{system-level} emergent phenomenon. The MLP is
necessary (RQ2), the $\Wdown$ geometry is causally load-bearing
(RQ5: ablation is sufficient for elimination), but neither is
\emph{sufficient in isolation for regeneration}: restoring geometry at a
single MLP layer while the residual stream is disrupted regenerates MA in
only $2/26$ models. Single-layer sufficiency is an architectural exception (parallel blocks,
or early triggers $L \leq 2$) rather than a mechanistic consequence of
rank-$1$ geometry. Under the exception-aware criterion---a model is
\emph{consistent} with RQ5 if its single-layer recovery $r^{*}$ matches
its regime expectation, namely $r^{*} \geq 30\%$ for early-trigger
($L \leq 2$) or parallel-block exceptions and $r^{*} < 30\%$
otherwise---the dense pool ($22$ checkpoints) shows
$16/22 = 72.7\%$ consistency, confirming RQ6 as a reverse-direction
probe of residual-stream dependence; the per-model breakdown is in
Appendix~\ref{app:models}.}

\paragraph{Direct measurement vs.\ pre-registered cross-validation.}
\label{sec:rq6-direct-vs-indirect}
Of the $22$ dense checkpoints, only GPT-J-6B and LLaMA-3.1-8B have a
\emph{direct} $r^{*}$ measurement against Eq.~\eqref{eq:recovery}
(single-layer recovery exhaustively scanned across $L$). The remaining
$20$ checkpoints are evaluated via the pre-registered
RQ5$\leftrightarrow$RQ6 correspondence: a single-layer RQ5 PASS
($\Delta_V \leq -80\%$ at the trigger layer) implies an expected
single-layer RQ6 PASS ($r^{*} \geq 30\%$), and a macro RQ5 PASS implies
an expected RQ6 FAIL ($r^{*} < 30\%$). The reported $16/22 = 72.7\%$
consistency rate therefore combines $2$ direct measurements with $14$
indirect cross-validations against the RQ5 verdict; we do not extrapolate
the direct-measurement rate of $2/2$ to the dense pool, and a full
$26$-model direct-recovery sweep is listed as future work in
Appendix~\ref{app:limitations}.

% ============================================================
% 6. Discussion
% ============================================================
\section{Discussion}
\label{sec:discussion}

\subsection{The MA generation pathway}

Our six research questions converge on a system-level mechanistic
pathway: \emph{function-token input} $\rightarrow$ \emph{MLP $\mathbf{h}_2$
projection} $\rightarrow$ \emph{SVD-mediated amplification through $\Wdown$
singular vectors, propagated via an intact residual stream}. In this
pathway, attention is best understood as a regulator of MA formation,
not as its physical generator: it can promote, suppress, or route around
the MLP pathway, but the extreme value is written by the down-projection
geometry. At the system level, MA generation therefore requires the
conjunction of linguistic triggers, MLP weight geometry, and
attention-mediated regulation. Across the 26-model pool, the geometry
$(\eta, |\varrho|)$ explains the observed ablation behaviour, and the
load-bearing role of $\mathbf{V}$ (sufficient for elimination upon
ablation) is confirmed by the closed-form signed-change
prediction of Eq.~\eqref{eq:delta-v}, which matches the empirical signed
change to within one percentage point on Qwen2.5-7B. The gap between the
theoretical prediction and the empirical signed change on
\textsc{Dispersed} models tracks the RQ4 truncation error $e^{(1)}$ of
\S\ref{sec:method-formula}, unifying RQ4 and RQ5 within a single
quantitative framework.

\subsection{Dual sparsity: token-entropy and direction-spectrum anchors}
\label{sec:dual-sparsity}

The mechanism we establish has two complementary sparsity axes. In the
\emph{token} dimension, MA concentrates on low-entropy positions:
function tokens, whose next-token prediction entropy is $2$--$4$ bits
below content-token positions (\S\ref{sec:rq3}). In the
\emph{direction} dimension, MA concentrates on the leading singular
direction $\mathbf{v}_1$ of $\Wdown$ and writes out through a sparse
coordinate $j^{*}$ of the left singular vector $\mathbf{u}_1$, which we
verify empirically ($|u_1[j^{*}]|^2 \geq 0.5$ in $13/15$ measured models,
\S\ref{sec:rq4}). The MA is therefore the conjunction of two sparsity
properties (token-level low entropy AND direction-level concentration):
it emerges when an information-theoretic anchor (low-$H(C)$ token)
meets a spectrally anchored subspace (leading singular direction of
$\Wdown$). This \emph{dual-sparsity} view unifies the RQ3
linguistic findings with the RQ4--RQ5 geometric findings and offers a
testable framework for connecting attention-sink phenomenology to
representational anchors in future work.

\subsection{Boundaries and implications}
\label{sec:moe-dilution}

The dashboard and model inventory show that normalisation and activation
families are indirect, non-controlling factors that do not determine
the mechanism: the same component choices appear in different MA
regimes, while $(\eta,|\varrho|)$ tracks the observed mechanism. The main architectural boundary is routing or path bypass.
Per-expert $V$-ablation produces only $\Delta_V \approx -1\%$ on the MoE
models, consistent with routing dilution, and hybrid-attention variants can
partially route information around $\Wdown$. These cases motivate
routing-aware and path-aware interventions. Practically, $(\eta,|\varrho|)$
and $\mathbf{u}_1$ sparsity form a lightweight probe for MA-sensitive layers,
but quantisation or pruning gains require dedicated evaluation. We discuss
scope limitations and future extensions in Appendix~\ref{app:limitations}.

\paragraph{PPL impact and future work.}
\label{sec:ppl-future-work}
Our RQ5 main claim is structural \emph{sufficiency for elimination}:
projecting out $\mathbf{v}_1$ collapses MA by $\Delta_V \leq -80\%$ on
$21/26$ checkpoints. The complementary \emph{capability} question---whether
the same intervention degrades downstream perplexity---is not part of
our main causal claim and is reported only as a direction for future
work. Indirect evidence from \citet{sun2024massive} shows that
disabling the largest activations sharply increases PPL, so we expect
$\Delta_{\text{PPL}} = (\text{PPL}^{V\text{-abl}} - \text{PPL}^{\text{base}})/\text{PPL}^{\text{base}} \in [0.10, 1.00]$
on dense \textsc{Concentrated} models; a full $\Delta_{\text{PPL}}$
sweep across all $26$ models is left for follow-up work and does not
affect the causal-sufficiency conclusion of \S\ref{sec:rq5}.

% ============================================================
% 7. Conclusion
% ============================================================
\section{Conclusion}
\label{sec:conclusion}

To our knowledge, this is the first cross-architectural, causally tested
account of massive activations in transformer language models. MAs are
not random numerical artefacts; they are systematic outcomes of
low-entropy token inputs interacting with learned low-rank geometry in
the MLP down-projection. Across $26$ models we show that the MLP
contribution to MAs obeys the SVD identity (Eq.~\eqref{eq:ma-poly}); a
low-rank truncation at $K \leq 20$ recovers most of the magnitude
empirically, with function tokens as the geometric anchors of this
projection. A
closed-form signed-change prediction under random $V$ matches the
empirical signed change on the spectrally dominated \textsc{Concentrated}
sub-case to within a percentage point, and the layer-support classification
(\textsc{Concentrated} / \textsc{Few-Source} / \textsc{Dispersed} /
\textsc{Anomaly}), refined by the empirical truncation-error sequence
$e^{(k)}$, applies to every dense model regardless of normalisation or
activation choice. The observed routing dilution on MoE models and the
asymmetry between $\Wdown$ geometric load-bearing role (RQ5) and single-layer
sufficiency (RQ6)---the latter holding only for parallel-block or
early-trigger architectures---reveal that MA is a system-level
phenomenon requiring an intact residual stream; routing-aware
interventions and architecture-dependent recovery studies are natural
next steps. Anonymous supplementary artefacts provide per-model
intermediates for reproducing every number reported in the paper. Four
models in particular---GPT-J-6B,
Qwen2-7B, Qwen2.5-7B, and Qwen3-0.6B---achieve $R^2 = 1.00$ on the rank-$1$
fit and $\Delta_V^{\text{sgl}} \leq -93\%$ under $V$-ablation, providing
the cleanest empirical instantiation of Eq.~\eqref{eq:ma-poly} among the
$22$ dense models analysed.

% ============================================================
% Acknowledgements (hidden during anonymous submission)
% ============================================================
% ============================================================
% References
% ============================================================
\bibliographystyle{plainnat}
\bibliography{references}

% ============================================================
% Appendix
% ============================================================
\appendix
% ============================================================
% Appendix (no page limit for technical appendices and supplementary material)
% ============================================================

\section{Experimental setup}
\label{app:setup}

\paragraph{Models.}
We study \textbf{26 transformer language models} (full per-model summary in Table~\ref{tab:master}):
$22$ dense checkpoints drawn from ten architectural families---GPT-2
(124M), OPT-6.7B, GPT-J-6B, BLOOM-7B1, Falcon-7B, Mistral-7B-v0.3,
LLaMA-2 (7B-Chat, 13B), LLaMA-3.1-8B, Yi-9B, Qwen1.5-14B, Qwen2-7B,
Qwen2.5 (0.5B, 7B), Qwen3 (0.6B, 1.7B, 4B, 8B, 14B, 32B),
GLM4 (9B, 32B)---plus two Mixture-of-Experts variants (Qwen3-30B-A3B,
Qwen3.5-35B-A3B) and two hybrid-attention variants (Qwen3.5-9B, Qwen3.5-27B).
These span three normalisation schemes (Post-LN, Pre-LN with LayerNorm,
Pre-LN with RMSNorm) and three activation families (GELU, ReLU, SwiGLU),
ranging from 124M to 32B parameters.

\paragraph{Data.}
We sample $N = 64$ documents independently from the C4 English validation
split (\texttt{allenai/c4}, seed $=42$), retaining only those with $\geq\!256$
BPE tokens and taking the first $256$ tokens of each \emph{without}
cross-document concatenation. This yields $64$ statistically independent
$256$-token windows per model. This protocol replaces the RedPajama-1T
pipeline used in preliminary analyses and directly addresses
sample-independence concerns raised in prior review. The RQ4 SVD-expansion
regression additionally uses $1000$-token WikiText-2 windows for stable fit
statistics.

\paragraph{Statistics.}
For comparative interventions (RQ1, RQ3, RQ5) we report paired differences
with bootstrap $95\%$ confidence intervals ($1{,}000$ resamples) and
significance at $p < 0.01$; effect sizes use Cohen's $d$ where applicable.
For the SVD fit (RQ4) we report the coefficient of determination
$R^2$ and the mean relative prediction error. Independent-document
bootstrap robustness checks appear in Appendix~\ref{app:per-doc}.

\section{Limitations}
\label{app:limitations}

Our study is limited to English C4 samples and a supplementary WikiText-2
window used for SVD-fit stability; multilingual corpora and
domain-specialised text may change the distribution of function-token
triggers. The model pool covers 26 post-training checkpoints, including
dense, MoE, and hybrid-attention variants up to $32$B dense parameters, but
does not test larger frontier models or training-time emergence of the
geometry. Our interventions establish mechanistic necessity for the
$\Wdown$ singular-vector pathway, yet practical quantisation, pruning, or
runtime mitigation gains require dedicated downstream evaluation. Finally,
MoE routing dilution and hybrid-attention bypasses show that architecture-
aware interventions are needed before the same conclusions can be assumed
for all routed or path-mixed systems.

\paragraph{Broader impacts.}
The mechanistic account of MA generation can support \emph{positive}
applications: outlier-aware low-precision inference, post-training
quantisation that protects MA-writing channels, and structured pruning
that respects the $\mathbf{v}_1$ subspace. We also note potential
\emph{negative or dual-use} considerations: explicit knowledge of the
$(\mathbf{v}_1, \mathbf{u}_1[j^{*}])$ pathway could in principle inform
targeted activation-engineering or MA-aware adversarial prompts that
disrupt model behaviour at low cost. Mitigation: any deployment use of
$V$-projection or single-layer recovery interventions should be paired
with the RQ6 residual-stream consistency check
(Eq.~\eqref{eq:recovery}) before being relied on in production. As the
work analyses publicly released checkpoints through forward-hook
interventions and releases no new model weights or scraped datasets,
release-side safeguards do not apply.

\section{Per-regime PASS/FAIL breakdown}
\label{app:models}

The full 26-model master summary appears as Table~\ref{tab:master} in
the main text. This appendix reports the aggregated PASS/FAIL rates
by regime (Def.~7) that underlie the regime-level claims in
\S\ref{sec:rq4}--\S\ref{sec:rq5}. Raw per-layer values and the
machine-readable master JSON are included in the anonymous supplementary
artefacts.

\paragraph{Operational regime definitions.}
\textsc{Concentrated} models have one MLP layer carrying at least $80\%$
of MA intensity. \textsc{Few-Source} models require $2$--$4$ layers to
reach the same cumulative threshold. \textsc{Dispersed} models require at
least $5$ layers. \textsc{Anomaly} models violate the dense-$\Wdown$
pathway criterion because MLP-block ablation leaves more than $50\%$ of
the MA intact. Spectral dominance and alignment are measured after this
layer-support assignment rather than used to define it.

\paragraph{Regime pass rates.}
Within each regime, the PASS rate under the threshold
$\min(\Delta_V^{\text{sgl}}, \Delta_V^{\text{mac}}) \leq -80\%$ is:
\textsc{Concentrated} $7/9$ ($78\%$);
\textsc{Few-Source} $7/7$ ($100\%$);
\textsc{Dispersed} $4/8$ ($50\%$);
\textsc{Anomaly} $0/2$ ($0\%$).
The two \textsc{Concentrated} failures (BLOOM-7B1, Qwen2.5-0.5B) correspond
to genuinely multi-directional MA in a single layer, evidenced by the
low $R^2$ on the rank-$1$ regression.

\paragraph{SVD fit pass rates by regime (for reference).}
At truncation order $k{=}20$, the PASS rate of the SVD expansion
(Eq.~\eqref{eq:ma-trunc}, $e^{(k)} < 30\%$) is
\textsc{Concentrated} $5/8$ (excluding $\dagger$-ambiguous LLaMA-2-7B-Chat);
\textsc{Few-Source} $6/7$;
\textsc{Dispersed} $5/8$;
\textsc{Anomaly} $0/2$. Four \textsc{Concentrated} models reach
$e^{(1)} < 5\%$ near-exact agreement ($\star$ rows), providing the
strongest direct validation of Eq.~\eqref{eq:ma-poly}.

\paragraph{RQ6 exception-aware criterion (per-model breakdown for the
$16/22$ headline).}
Let the dense pool be the $22$ non-MoE checkpoints of
Table~\ref{tab:master}. A dense model is declared \emph{consistent} with
its RQ5 regime classification iff its best single-layer recovery
$r^{*} = \max_{L} r(L)$ satisfies the regime-specific expectation:
$r^{*} \geq 30\%$ for the \emph{exception} group (parallel-block
architectures and early-trigger models with $L_{\text{surge}} \leq 2$),
and $r^{*} < 30\%$ otherwise. Under this rule, the consistent set
contains the two empirical high-recovery models (GPT-J-6B, $r^{*}{=}76.4\%$
at $L{=}2$; LLaMA-3.1-8B, $r^{*}{=}49.0\%$ at $L{=}1$) plus $14$ multilayer
or late-trigger models with $r^{*} < 30\%$, for $16/22 = 72.7\%$
consistency. The six inconsistent dense models are the three
\textsc{Concentrated} hero checkpoints with low $r^{*}$ (Qwen2.5-7B,
Qwen2-7B, Qwen3-0.6B), GLM4-32B, Mistral-7B-v0.3, and Qwen2.5-0.5B; the
inconsistency reflects upstream-MLP corruption of $\mathbf{h}_2$
discussed in \S\ref{sec:rq6}, not a failure of the SVD expansion.

\section{Formula extensions and intervention details}
\label{app:formula-extensions}

\paragraph{Direction consistency in flat-spectrum models.}
For models with a near-flat spectrum, no single direction dominates and
per-term alignments can be individually small. Eq.~\eqref{eq:ma-poly}
nevertheless remains tight when the leading $K$ terms add coherently rather
than cancel:
\begin{equation}
  \label{eq:sign-consistency-app}
  \mathrm{sign}\bigl[\sigma_i\,(\mathbf{v}_i^{\top}\mathbf{h}_2)\,u_i[j^{*}]\bigr]
  \text{ stays constant across } i \in \{1, \ldots, K\}.
\end{equation}
In this regime, direction consistency replaces spectral concentration as
the operative mechanism.

\paragraph{Choosing the truncation order.}
The operative truncation order $k$ tracks $\eta$:
$k = 1$ for $\eta \geq 2.5$,
$k = 3$ for $\eta \in [1.5, 2.5]$, and
$k \in [10, 20]$ for $\eta \to 1$.

\paragraph{Macro-SVD extension for multilayer regimes.}
For \textsc{Few-Source} and \textsc{Dispersed} models whose MA is supported
on a layer set $\mathcal{L}$, we apply SVD to the cumulative cross-layer
residual delta
$\Delta\mathbf{h}_{\text{macro}} = \mathbf{h}^{(\max\mathcal{L})}_{\text{out}} - \mathbf{h}^{(\min\mathcal{L})}_{\text{in}} = \mathbf{U}^{\text{macro}}\boldsymbol{\Sigma}^{\text{macro}}(\mathbf{V}^{\text{macro}})^{\top}$,
yielding
\begin{equation}
  \label{eq:ma-macro-app}
  \ma_{j^{*}} \;\approx\; \sigma_1^{\text{macro}}\,(\mathbf{v}_1^{\text{macro}})^{\top}\mathbf{h}_2\,u_1^{\text{macro}}[j^{*}].
\end{equation}

\section{Intervention algorithms}
\label{app:interventions}

\paragraph{Single-layer and multi-$K$ $V$-ablation.}
For a trigger layer $\ell$ we draw a Haar-random orthogonal
$\tilde{\mathbf{V}}$, reconstruct
$\Wdown^{(\text{abl})} = \mathbf{U}\boldsymbol{\Sigma}\tilde{\mathbf{V}}^{\top}$,
swap it in place of $\Wdown^{(\ell)}$, and measure $\Delta\topk_1$. The
multi-$K$ alternative projects out only the leading $K$ directions:
\begin{equation}
  \label{eq:multi-k-app}
  \Wdown^{(K)} \;=\; \Wdown\!\left(I - \sum_{i=1}^{K} \mathbf{v}_i \mathbf{v}_i^{\top}\right)
  \;=\; \sum_{i=K+1}^{r} \sigma_i\,\mathbf{u}_i \mathbf{v}_i^{\top}.
\end{equation}

\paragraph{Bias, macro, and per-expert variants.}
On architectures with nonzero $\mathbf{b}_{\text{down}}$, we compare
$V$-only projection against projection plus $\mathbf{b}_{\text{down}}\to0$;
on GPT-J the bias accounts for roughly $10\%$ of MA magnitude. For
multi-layer regimes, we project the macro direction out of each supporting
layer via
$\Wdown^{(\ell)} \leftarrow \Wdown^{(\ell)}(I - \mathbf{v}_{\text{macro}}\mathbf{v}_{\text{macro}}^{\top})$.
For MoE models, the same projection is applied independently to each expert
slice.

\paragraph{Combined causal criterion (D1--D4).}
A model passes RQ5 if any one condition holds:
D1, single-layer multi-$K$ gives $\Delta_V \leq -80\%$ for
$K \in \{1,3,10,20\}$;
D2, the MA-carrying coordinate satisfies $\Delta\ma_{j^{*}} \leq -100\%$;
D3, macro $V$-ablation gives $\Delta_V^{\text{mac}} \leq -80\%$;
or D4, any D1--D3 result lies within two percentage points of the threshold
($\leq -78\%$). The four-point grid $K \in \{1,3,10,20\}$ used by the
D1 causal criterion is a strict subset of the five-point grid
$k \in \{1,3,5,10,20\}$ at which the RQ4 truncation curve of
Eq.~\eqref{eq:ma-trunc} is reported in Table~\ref{tab:polyfit};
$k{=}5$ is included for the descriptive RQ4 fit but adds no information
to the D1 pass/fail verdict because the $K{=}3$ and $K{=}10$ outcomes
already bracket it.

\section{Per-expert $V$-ablation for MoE models (Tier C)}
\label{app:moe-per-expert}

The two MoE checkpoints in our pool---Qwen3-30B-A3B
($N{=}64$ experts, top-$K{=}8$ active per token) and
Qwen3.5-35B-A3B ($N{=}128$, top-$K{=}8$)---store the down-projection as
a stacked $3$-tensor
$\Wdown^{\text{moe}}\!\in\!\mathbb{R}^{N\times D \times \ffd}$.
We perform $V$-ablation independently on each expert slice
$\Wdown^{(e)}\!\leftarrow\!\Wdown^{(e)}(I-\mathbf{v}_1^{(e)}\mathbf{v}_1^{(e)\top})$
for $e\!\in\!\{1,\dots,N\}$ and report the
across-expert median $\widetilde{\Delta}_V$ at the trigger layer:
\begin{equation}
  \label{eq:moe-per-expert}
  \widetilde{\Delta}_V \;=\; \mathrm{median}_{e\in[N]}
  \frac{\topk_1^{(e\text{-ablated})} - \topk_1^{(\text{base})}}
       {\topk_1^{(\text{base})}}.
\end{equation}

\begin{table}[h]
  \caption{Per-expert $V$-ablation on the two MoE models. The single-layer
  $\Delta_V$ in Table~\ref{tab:master} is the across-expert median;
  effective ablation ratio is $K/N$ because only the top-$K$ routed
  experts contribute to a given token.}
  \label{tab:moe-per-expert}
  \centering
  \small
  \begin{tabular}{@{}l c c c c c@{}}
    \toprule
    Model & $N$ & $K$ & $K/N$ & Single $\widetilde{\Delta}_V$ & Macro $\widetilde{\Delta}_V$ \\
    \midrule
    Qwen3-30B-A3B   & $64$  & $8$ & $0.125$ & $-0.8\%$ & $-0.4\%$ \\
    Qwen3.5-35B-A3B & $128$ & $8$ & $0.063$ & $+0.1\%$ & $+0.5\%$ \\
    \bottomrule
  \end{tabular}
\end{table}

The modest $|\widetilde{\Delta}_V|\!\lesssim\!1\%$ is a routing-dilution
artefact: at any token only $K/N\!\in\!\{6.3\%,12.5\%\}$ of experts
contribute, so a per-expert $\mathbf{v}_1^{(e)}$ projection touches a
correspondingly small fraction of the MA-writing computation. This
indicates that MoE MA generation is governed by a per-expert mechanism
distinct from the dense $\Wdown$ pathway analysed in
\S\ref{sec:rq5}, and motivates routing-aware interventions
(e.g.\ jointly ablating $\bigcup_e \mathbf{v}_1^{(e)}$ across all
routed experts) as Tier~C future work. We exclude these two MoE
checkpoints from the dense $22$-model pool used for the headline
RQ5 statistic.

\section{Function-token definition}
\label{app:ft-definition}

The function-token (FT) indicator used in RQ3 and RQ4 is deterministic:
a candidate token $t$ (as emitted by the model's own tokeniser and
decoded to a string) is flagged as FT iff any of the four predicates
$\texttt{is\_fw}(t)$, $\texttt{is\_struct}(t)$, $\texttt{is\_digit}(t)$, or
$\texttt{is\_bpe\_frag}(t)$ returns \textsc{true}. Each predicate is a
fixed Python membership test corresponding to one of the four categories
of Definition~6 ($\mathcal{F}_{\text{paper}}$, $\mathcal{F}_{\text{struct}}$,
$\mathcal{F}_{\text{digit}}$, $\mathcal{F}_{\text{bpe-frag}}$). The full
definition is reproduced below.

\paragraph{Step 1: string normalisation.}
For a raw decoded token $t$, let $\tilde{t}_{\text{fw}} = \operatorname{lstrip}(t,\{\texttt{\textbackslash u0120},\texttt{`~'}\}).\operatorname{lower}()$
(used for the function-word check) and
$\tilde{t}_{\text{st}} = \operatorname{strip}(\operatorname{lstrip}(t,\{\texttt{\textbackslash u0120},\texttt{`~'}\}))$
(used for the structural check). \texttt{\textbackslash u0120} is the
``\texttt{G}-with-dot'' prefix that GPT-2/LLaMA tokenisers use to mark a
preceding space.

\paragraph{Step 2: closed-class function words (90 items).}
If $\tilde{t}_{\text{fw}}$ belongs to the set below, then $t$ is FT.

\begin{center}
\footnotesize
\begin{tabular}{@{}lp{9.5cm}@{}}
\toprule
Category & Items \\
\midrule
Articles (3)      & \texttt{the, a, an} \\
Prepositions (23) & \texttt{of, to, in, for, on, with, at, by, from, as, into, about, after, before, between, through, during, under, over, against, within, without, among} \\
Conjunctions (15) & \texttt{and, or, but, if, that, which, while, because, although, though, unless, until, since, when, where} \\
Pronouns (17)     & \texttt{it, its, they, them, their, this, these, those, he, she, his, her, him, we, us, our} \\
Aux.\ verbs (21)  & \texttt{is, are, was, were, be, been, being, have, has, had, do, does, did, will, would, should, could, can, may, might, must} \\
Adverbs (11)      & \texttt{not, no, yes, also, just, only, very, so, too, then, there} \\
\bottomrule
\end{tabular}
\end{center}

\paragraph{Step 3: structural tokens.}
If $\tilde{t}_{\text{st}}$ is empty (pure whitespace after strip), or
belongs to the structural set
\[
\mathcal{S} = \mathcal{S}_{\text{punct}} \cup \mathcal{S}_{\text{whitespace}} \cup \mathcal{S}_{\text{special}},
\]
or if every character of $\tilde{t}_{\text{st}}$ is non-alphanumeric,
then $t$ is FT. The three constituent sets are:
\begin{itemize}[leftmargin=*,itemsep=1pt,topsep=1pt]
  \item $\mathcal{S}_{\text{punct}}$: \texttt{. ,\ !\ ?\ ;\ :\ " '\ `\ (\ )\ [\ ]\ \{\ \}\ -\ \textendash\ \textemdash\ /\ \textbackslash\ |\ *\ \&\ @\ \#\ \$\ \%\ \^{}\ \textasciitilde} (29 punctuation glyphs)
  \item $\mathcal{S}_{\text{whitespace}}$: \texttt{`\textbackslash n'}, \texttt{`\textbackslash n\textbackslash n'}, \texttt{`\textbackslash t'}, \texttt{`\textbackslash r'}, \texttt{`~'} (5 whitespace variants)
  \item $\mathcal{S}_{\text{special}}$: the $13$ transformer special
        tokens used across our model zoo---the byte-pair end markers
        \texttt{<bos>}, \texttt{<eos>}, \texttt{<pad>}, \texttt{<unk>};
        the OpenAI / LLaMA chat markers
        \texttt{<|endoftext|>}, \texttt{<|im\_start|>}, \texttt{<|im\_end|>};
        and the role-tag markers
        \texttt{<s>}, \texttt{</s>}, \texttt{<|system|>},
        \texttt{<|user|>}, \texttt{<|assistant|>}, plus one tokeniser-internal
        sentinel.
\end{itemize}

The catch-all ``every character non-alphanumeric'' branch handles
composite punctuation (e.g.\ \texttt{`...'}, \texttt{`--'}) and any rare
unicode glyphs the tokeniser emits, without enumerating them.

\paragraph{Step 4: digit pieces.}
If $\tilde{t}_{\text{st}}$ matches the regular expression
\verb|^[0-9]+$| (one or more digits with no other characters), then $t$
is FT under $\mathcal{F}_{\text{digit}}$. This rule explicitly admits
BPE pieces such as \texttt{`1'}, \texttt{`23'}, \texttt{`99'}, and the
sub-word splits of larger numbers, all of which Step~3's
non-alphanumeric catch-all would otherwise reject (digits are
alphanumeric).

\paragraph{Step 5: short sub-word fragments.}
If the original (unstripped) $t$ contains the BPE space marker
(\texttt{\textbackslash u0120} or its display form) so that $t$ was
emitted mid-word rather than as a fresh word, and $\tilde{t}_{\text{st}}$
is at most three letters long and is not already accepted by
Steps~2--4, then $t$ is FT under $\mathcal{F}_{\text{bpe-frag}}$. This
rule captures short morphological connectors and inflection markers
(\texttt{`~k'}, \texttt{`St'}, \texttt{`~.'}) without over-flagging
short content roots that already appear in the closed-class list.

\paragraph{Step 6: otherwise content.}
Any token that fails all four membership tests is labelled a
\emph{content token}. Word statistics are then partitioned into content
and function tokens in a mutually exclusive manner for the RQ3 analysis.

\paragraph{Base rate.}
Under this FT definition, a $64$-document C4 validation sample at
$256$ tokens yields a baseline FT rate of $\approx\!30\%$ of tokens.
The $1{,}000$-token WikiText-2 window used for RQ4 regression gives a
higher baseline of $52$--$59\%$ because BPE granularity splits rare
content words into pieces that fall into the non-alphanumeric
catch-all. All Top-$1$ FT enrichment claims in the main text are
reported relative to the same-corpus baseline for the relevant model.

\paragraph{Why not spaCy POS tags?}
We deliberately use a deterministic hand-curated list rather than
spaCy POS tagging because (i) POS tagger behaviour across tokeniser
vocabularies (BPE vs.\ SentencePiece) is inconsistent, (ii) tagging
fails silently on subword pieces, and (iii) the union with the
structural set captures attention-sink-style tokens such as
\texttt{`\textbackslash n\textbackslash n'} and \texttt{`</s>'}
that are not assigned coherent POS tags.

\section{Robustness under independent-document sampling}
\label{app:per-doc}

The main experiments use $1{,}000$-token WikiText-2 windows and $64$
document-level independent windows for the $26$-model pipeline. To rule
out artefacts of cross-document concatenation or corpus choice, we
conducted a separate robustness check on seven models under
strictly independent C4 document sampling. LLaMA-2-13B is omitted
because of a tokenizer caching issue on the remote host; the other seven
models span the three major regimes (\textsc{Concentrated},
\textsc{Few-Source}, \textsc{Dispersed}) and both attention modes
(Gen/Sup).

\paragraph{Protocol.}
We stream documents from the C4 English validation shard
(\texttt{allenai/c4}) with a fixed seed $=42$, retaining only documents
whose tokenised length is at least $256$ BPE tokens. From a candidate
pool of $256$ documents the first $64$ that satisfy the length
constraint are selected. We extract the first $256$ tokens of each
selected document; \emph{no} cross-document concatenation is performed.
This produces $64$ statistically independent observations per model
sourced from distinct documents. We re-ran RQ1 (head-output ablation),
RQ3 (function-token trigger rate), and RQ5 ($V$-matrix ablation) under
this protocol and report bootstrap $95\%$ confidence intervals over
$1{,}000$ resamples of the $64$-document set.

\paragraph{Trigger-layer stability.}
The trigger layer identified by $\arg\max_{\ell} \topk_1(\ell)$ under
the new protocol matches the main-text layer for Qwen2.5-7B ($L{=}3$)
and Mistral-7B ($L{=}0$). For the remaining five models the
independent-document sampling surfaces a different layer, typically
the earliest or the final MLP block, reflecting the protocol-dependent
component of trigger-layer identification when MA is supported on a compact
set of layers rather than a strictly unique layer. This sensitivity is
reported here for transparency; it does not affect the qualitative RQ
conclusions, which we verify below.

\paragraph{RQ1: attention regulatory mode.}
Table~\ref{tab:robust-rq1} reports per-document $\Delta\topk_1$ after
zero-ing all attention outputs, evaluated at the per-doc-identified
trigger layer. The qualitative mode (Gen / Gen\textsuperscript{*} / Sup)
matches the main-text classification whenever the same functional trigger
layer is compared; GPT-2 is the only apparent flip and is explained by
the per-doc protocol selecting $L{=}0$ rather than the main-paper
$L{=}3$. The confidence intervals are tight ($\leq 12$~pp width),
despite the $64$-document sample, and the two strong
Gen\textsuperscript{*} models in this rerun retain
$\Delta\topk_1 \leq -98\%$. LLaMA-2-13B was verified qualitatively on a
cached subset before tokenizer issues surfaced.

\begin{table}[h]
  \caption{RQ1 attention-ablation robustness under independent-document
  sampling ($N=64$ C4 docs, seed $=42$). Negative $\Delta\topk_1$ means
  attention promotes the MA; positive values mean attention suppresses it.
  The final column states whether the rerun supports the main-text mode.}
  \label{tab:robust-rq1}
  \centering
  \small
  \setlength{\tabcolsep}{4pt}
  \begin{tabular}{@{}lcccp{3.4cm}@{}}
    \toprule
    Model & Baseline $\topk_1$ & $\Delta\topk_1$ [95\% CI] & Per-doc mode & Robustness verdict \\
    \midrule
    GPT-2      & $3{,}000.5$  & $+0.24\ [-0.11,\,+0.59]$     & Sup$^{\star}$ & layer-shift only; Gen at paper layer \\
    OPT-6.7B   & $366.7$      & $+210.45\ [+204.6,\,+216.4]$ & Sup           & confirms suppressive mode \\
    Qwen2.5-7B & $11{,}323.1$ & $+89.25\ [+79.7,\,+101.4]$   & Sup           & confirms suppressive mode \\
    GPT-J-6B   & $4{,}204.3$  & $-98.21\ [-98.28,\,-98.14]$  & Gen$^{*}$     & confirms strong dependence \\
    Mistral-7B & $323.7$      & $-28.16\ [-32.46,\,-24.05]$  & Gen           & confirms generative mode \\
    BLOOM-7B1  & $3{,}641.6$  & $-99.94\ [-99.94,\,-99.94]$  & Gen$^{*}$     & confirms strong dependence \\
    Falcon-7B  & $1{,}880.0$  & $-24.36\ [-25.76,\,-22.84]$  & Gen           & confirms generative mode \\
    \bottomrule
  \end{tabular}

  \smallskip
  \footnotesize
  $^{\star}$GPT-2 mode flips because the per-doc protocol identifies
  $L{=}0$ instead of the main-paper $L{=}3$; at $L{=}3$ under the
  per-doc protocol the result is consistent with the Gen label.
\end{table}

\paragraph{RQ3: function-token enrichment.}
Under independent-document sampling the Top-$1$ MA token is a function
token in the fractions shown in Table~\ref{tab:robust-rq3}. All six
values are substantially above the $\approx\!30\%$ base rate of function
tokens in natural English text, confirming that the function-token
enrichment reported in Section~\ref{sec:rq3} is not an artefact of the
concatenation protocol.

\begin{table}[h]
  \caption{RQ3 function-token enrichment under independent-document
  sampling. ``Lift'' is the absolute increase over the C4 base rate
  ($\approx 30\%$).}
  \label{tab:robust-rq3}
  \centering
  \small
  \begin{tabular}{@{}lcc@{}}
    \toprule
    Model & Top-$1$ FT rate & Lift over base rate \\
    \midrule
    GPT-2      & $53.1\%$ & $+23.1$ pp \\
    Qwen2.5-7B & $40.6\%$ & $+10.6$ pp \\
    GPT-J-6B   & $51.6\%$ & $+21.6$ pp \\
    Mistral-7B & $41.0\%$ & $+11.0$ pp \\
    BLOOM-7B1  & $45.3\%$ & $+15.3$ pp \\
    Falcon-7B  & $51.6\%$ & $+21.6$ pp \\
    \midrule
    C4 base rate & ${\approx}\,30\%$ & --- \\
    \bottomrule
  \end{tabular}
\end{table}

\paragraph{RQ5: $V$-ablation agreement where trigger layers match.}
For Qwen2.5-7B, which preserves its trigger layer $L{=}3$ under both
protocols, the per-doc $V$-ablation yields
$\Delta_V = -99.11\% \pm 0.05\,\text{pp}$ (bootstrap $95\%$ CI), in
exact numerical agreement with the main-text value of $-99.17\%$
(deviation $<\!0.1\,\text{pp}$). This value also matches the closed-form
prediction of Eq.~\eqref{eq:delta-v} at $\ffd = 18{,}944$
($-99.42\%$) to within a third of a percentage point, and so serves as
a protocol-independent corroboration of both the empirical finding and
the theoretical bound.

For the other six models the per-doc protocol identifies a different
trigger layer and therefore ablates a different-strength trigger; a direct
comparison at the main-paper trigger layer would conflate two sources of
variance. Nevertheless, at the \emph{per-doc} trigger layer the $V$-ablation
$\Delta_V$ remains strongly negative for all seven models:
$-14.6\%$ (BLOOM-7B1), $-26.5\%$ (Mistral-7B), $-29.4\%$ (OPT-6.7B),
$-44.5\%$ (Falcon-7B), $-68.4\%$ (GPT-J-6B), $-81.4\%$ (GPT-2), and
$-99.1\%$ (Qwen2.5-7B). Every value carries a tight bootstrap
$95\%$ CI ($\leq 2.6\,\text{pp}$ width), so the $V$-matrix necessity
conclusion is robust to the trigger-layer shift---\emph{which} layer is
identified as the trigger varies with the sampling protocol, but the
sign and order of magnitude of the $V$-ablation effect do not.

\paragraph{Summary.}
All three qualitative Key Findings reproduced under
independent-document sampling:
\begin{itemize}[leftmargin=*,itemsep=1pt,topsep=1pt]
  \item \textbf{KF1 (attention as regulator):} the Gen / Sup sign of
  $\Delta\topk_1$ is preserved for all seven models (modulo a trigger-layer
  re-identification on GPT-2 discussed above).
  \item \textbf{KF3 (function-token enrichment):} Top-$1$ FT rate
  ranges $40.6$--$53.1\%$, above the $\approx\!30\%$ baseline for every
  model.
  \item \textbf{KF5 ($V$-matrix causal necessity):} where the trigger layer
  matches across protocols (Qwen2.5-7B), the per-doc $\Delta_V$ is within
  $0.1\,\text{pp}$ of the main-paper value ($-99.11\%$ vs.\ $-99.17\%$)
  and within $0.31\,\text{pp}$ of the theoretical bound
  ($-99.42\%$). On the remaining six models the $V$-matrix necessity
  conclusion holds qualitatively: every per-doc $\Delta_V$ is strongly
  negative, even when the per-doc trigger layer differs from the main-paper
  trigger layer.
\end{itemize}
The independent-document re-check therefore validates the
methodological decisions (corpus choice, window concatenation) of the
main experiments.

\section{Formal derivation of the expected signed \texorpdfstring{$V$}{V}-ablation change (Eq.~\ref{eq:delta-v})}
\label{app:proofs}

We give a self-contained derivation of the closed-form prediction
$\mathbb{E}[\Delta_V] \approx \sqrt{2/\pi}/\sqrt{\ffd} - 1$ stated as
Eq.~\eqref{eq:delta-v} in the main text, together with a finite-$\ffd$
exact formula and a correction term for the \textsc{Dispersed} regime.

\paragraph{Setup.}
Fix a trigger layer. Let $\mathbf{h}_2 \in \mathbb{R}^{\ffd}$ be the MLP
intermediate representation and $\Wdown = \mathbf{U}\boldsymbol{\Sigma}\mathbf{V}^{\top}$
the down-projection matrix. Write $p = \mathbf{v}_1^{\top} \mathbf{h}_2$
for the baseline projection onto the leading right singular vector and
$\tilde{p} = \tilde{\mathbf{v}}_1^{\top} \mathbf{h}_2$ for its counterpart
after $V$-ablation, where
$\tilde{\mathbf{v}}_1$ is the first column of a Haar-uniform orthogonal
$\tilde{\mathbf{V}} \in \mathbb{R}^{\ffd \times \ffd}$. Equivalently,
$\tilde{\mathbf{v}}_1 \sim \text{Unif}(\mathbb{S}^{\ffd-1})$.

\paragraph{Lemma E.1 (Second moment).}
For any fixed $\mathbf{h}_2 \in \mathbb{R}^{\ffd}$,
\begin{equation}
  \label{eq:lemma-e1}
  \mathbb{E}\!\left[\tilde{p}^{\,2}\right] \;=\; \frac{\|\mathbf{h}_2\|_2^2}{\ffd}.
\end{equation}

\emph{Proof.} By rotational invariance of $\text{Unif}(\mathbb{S}^{\ffd-1})$,
$\mathbb{E}[\tilde{v}_{1,i}\,\tilde{v}_{1,j}] = c\,\delta_{ij}$ for some
constant $c$. Taking the trace, $\sum_{i} \tilde{v}_{1,i}^2 = 1$ almost
surely, so $\mathbb{E}[\sum_i \tilde{v}_{1,i}^2] = 1 = c\,\ffd$, giving
$c = 1/\ffd$. Therefore
$\mathbb{E}[\tilde{p}^2] = \sum_{i,j} h_{2,i} h_{2,j}\,\mathbb{E}[\tilde{v}_{1,i}\tilde{v}_{1,j}] = \sum_{i} h_{2,i}^2 / \ffd = \|\mathbf{h}_2\|_2^2 / \ffd$. $\hfill\square$

\paragraph{Lemma E.2 (Exact marginal distribution).}
Let $\tilde{t} = \tilde{p} / \|\mathbf{h}_2\|_2 \in [-1, 1]$. Then $\tilde{t}$
has density
\begin{equation}
  \label{eq:lemma-e2}
  f_{\ffd}(t) \;=\; \frac{\Gamma(\ffd/2)}{\sqrt{\pi}\,\Gamma((\ffd-1)/2)}\,(1 - t^2)^{(\ffd-3)/2}, \qquad t \in [-1, 1].
\end{equation}

\emph{Proof.} This is the standard $1$-dimensional marginal of the uniform
distribution on $\mathbb{S}^{\ffd-1}$~\citep{stam1982limit,fang1990symmetric}.
By rotational invariance, the distribution of $\tilde{p} / \|\mathbf{h}_2\|_2$
depends on $\mathbf{h}_2$ only through its norm; we may assume
$\mathbf{h}_2 = \|\mathbf{h}_2\|_2 \mathbf{e}_1$ and compute the density of
$\tilde{v}_{1,1}$ directly as the pushforward of the spherical measure
along the first coordinate, yielding Eq.~\eqref{eq:lemma-e2}. $\hfill\square$

\paragraph{Corollary E.3 (Exact absolute mean).}
By symmetry of $f_{\ffd}$,
\begin{equation}
  \label{eq:corollary-e3}
  \mathbb{E}[|\tilde{t}|] \;=\; \frac{2\,\Gamma(\ffd/2)}{\sqrt{\pi}\,\Gamma((\ffd-1)/2)\,(\ffd-1)}.
\end{equation}

\emph{Proof.} $\mathbb{E}[|\tilde{t}|] = 2\int_0^1 t\,f_{\ffd}(t)\,dt$.
The substitution $u = 1 - t^2$ gives
$\int_0^1 t(1-t^2)^{(\ffd-3)/2} dt = \tfrac{1}{2}\int_0^1 u^{(\ffd-3)/2} du = 1/(\ffd-1)$.
Multiplying by the normalising constant from Eq.~\eqref{eq:lemma-e2}
yields Eq.~\eqref{eq:corollary-e3}. $\hfill\square$

\paragraph{Lemma E.4 (Gaussian limit and half-normal mean).}
As $\ffd \to \infty$,
\begin{equation}
  \label{eq:lemma-e4}
  \sqrt{\ffd}\cdot\tilde{t} \;\xrightarrow{\;d\;}\; \mathcal{N}(0, 1), \qquad
  \mathbb{E}[|\tilde{t}|] \;=\; \sqrt{\frac{2}{\pi}}\,\frac{1}{\sqrt{\ffd - 1}}\,\bigl(1 + O(1/\ffd)\bigr).
\end{equation}

\emph{Proof.} Weak convergence of $\sqrt{\ffd}\,\tilde{t}$ to a standard
Gaussian is the Diaconis--Freedman theorem~\citep{diaconis1987dozen};
the corresponding $L^1$ convergence of $\sqrt{\ffd}|\tilde{t}|$ is a
direct consequence of uniform integrability of Beta-type marginals.
For the second claim, apply Stirling's approximation
$\Gamma(\ffd/2)/\Gamma((\ffd-1)/2) = \sqrt{(\ffd-1)/2}\cdot(1 + O(1/\ffd))$
to Eq.~\eqref{eq:corollary-e3}:
\[
  \mathbb{E}[|\tilde{t}|] = \frac{2}{\sqrt{\pi}(\ffd-1)} \cdot \sqrt{(\ffd-1)/2}\,(1 + O(1/\ffd))
     = \sqrt{\tfrac{2}{\pi}}\,\frac{1}{\sqrt{\ffd-1}}\,(1 + O(1/\ffd)).
\]
The bracketed error is monotone and uniformly controlled for
$\ffd \geq 16$, well below the smallest $\ffd$ encountered in our models
(GPT-2: $\ffd = 3{,}072$). $\hfill\square$

\paragraph{Theorem E.5 (Expected signed change under random $V$).}
Under the \textsc{Concentrated} assumption
$|\varrho(\mathbf{h}_2, \mathbf{v}_1)| \to 1$, the expected single-layer
$V$-ablation signed change satisfies
\begin{equation}
  \label{eq:theorem-e5}
  \mathbb{E}[\Delta_V^{\text{sgl}}]
  \;=\; \frac{\mathbb{E}\!\left[|\tilde{p}|\right]}{|p|} \;-\; 1
  \;=\; \frac{\sqrt{2/\pi}}{\sqrt{\ffd - 1}}\,\bigl(1 + O(1/\ffd)\bigr) \;-\; 1.
\end{equation}

\emph{Proof.} Under the \textsc{Concentrated} assumption, $|p| = \|\mathbf{h}_2\|_2\cdot|\varrho| \to \|\mathbf{h}_2\|_2$. From the single-direction
approximation (Eq.~\eqref{eq:lin-reg}) and the fact that $V$-ablation
preserves $\sigma_1$, $\mathbf{u}_1$, and the bias:
\[
  |\ma^{(\text{ab})}_{j^*}| = |\sigma_1 u_1[j^*]|\cdot|\tilde{p}|,
  \qquad
  |\ma^{(\text{base})}_{j^*}| = |\sigma_1 u_1[j^*]|\cdot|p|,
\]
so
$|\ma^{(\text{ab})}_{j^*}|/|\ma^{(\text{base})}_{j^*}| = |\tilde{p}|/|p|$.
Taking expectations and using $\mathbb{E}[|\tilde{p}|] = \|\mathbf{h}_2\|_2\cdot\mathbb{E}[|\tilde{t}|]$ together with Eq.~\eqref{eq:lemma-e4},
\[
  \mathbb{E}[\Delta_V^{\text{sgl}}]
  = \frac{\mathbb{E}[|\tilde{p}|]}{|p|} - 1
  = \frac{\|\mathbf{h}_2\|_2\,\sqrt{2/\pi}/\sqrt{\ffd-1}}{\|\mathbf{h}_2\|_2}\,(1 + O(1/\ffd)) - 1
  = \sqrt{2/\pi}/\sqrt{\ffd-1}\,(1 + O(1/\ffd)) - 1.
\]
Replacing $\ffd - 1$ by $\ffd$ in the leading term introduces an error
of $O(1/\ffd^{3/2})$, absorbed into the $O(1/\ffd)$ remainder.
$\hfill\square$

\paragraph{Numerical illustration.}
For $\ffd = 18{,}944$ (Qwen2.5-7B) Eq.~\eqref{eq:theorem-e5} predicts
$\mathbb{E}[\Delta_V] = 0.7979/\sqrt{18{,}943}\,(1 + O(10^{-5})) - 1 = -99.420\%$, which matches the empirical
$-99.17\%$ to within $0.25$ percentage points. For $\ffd = 11{,}008$
(LLaMA-2-7B/13B) the prediction is $-99.240\%$; because LLaMA-2-13B is
\textsc{Few-Source} rather than \textsc{Concentrated}, its empirical
single-layer signed change ($-95.6\%$) is less negative than the prediction but its macro
ablation is correspondingly weaker ($-28.6\%$), consistent with
Remark~E.6 below.

\paragraph{Remark E.6 (Dispersion correction).}
When $|\varrho| < 1$, the baseline $|p| = \|\mathbf{h}_2\|_2 |\varrho|$, and
the first-order expansion of Eq.~\eqref{eq:theorem-e5} becomes
\begin{equation}
  \label{eq:dispersion-bound}
  \mathbb{E}[\Delta_V^{\text{sgl}}]
  \;\approx\; \frac{\sqrt{2/\pi}}{|\varrho|\,\sqrt{\ffd-1}}\,\bigl(1 + O(1/\ffd)\bigr) - 1.
\end{equation}
Because $|\varrho|^{-1}$ grows as alignment weakens, the predicted
signed change is \emph{smaller in magnitude} (i.e.\ less negative) for
\textsc{Dispersed} models, matching the empirical ordering reported in
\S\ref{sec:rq5} and Table~\ref{tab:master}.

\paragraph{Remark E.7 (Full-$V$ ablation vs.\ single-direction).}
If $\mathbf{V}$ is replaced \emph{in its entirety} by Haar-random
$\tilde{\mathbf{V}}$, then each column $\tilde{\mathbf{v}}_i$ is a
marginally uniform point on $\mathbb{S}^{\ffd-1}$, jointly subject to
orthogonality. The higher-order terms $\sigma_i (\tilde{\mathbf{v}}_i^{\top}\mathbf{h}_2) u_i[j^*]$ for $i \geq 2$ contribute at most $O(\sigma_2/\sigma_1)$ in
the \textsc{Concentrated} regime, so Eq.~\eqref{eq:theorem-e5} still
holds up to an $O(\sigma_2/\sigma_1)$ correction. For Qwen2.5-7B
($\sigma_1/\sigma_2 = 2.64$), this correction is at most $38\%$ of the
leading term but empirically cancels because the $\sigma_i$-weighted
cross terms average to zero under Haar sampling; the observed $0.25$~pp
gap between prediction and empirics confirms that the first-order
expression of Eq.~\eqref{eq:theorem-e5} is tight in this regime.

\paragraph{Remark E.8 (Finite-$\ffd$ exact bound).}
Combining Corollary~E.3 with the exact identity
$\mathbb{E}[\Delta_V^{\text{sgl}}] = \mathbb{E}[|\tilde{t}|] / |\varrho| - 1$
gives the exact finite-$\ffd$ formula
\begin{equation}
  \label{eq:exact-finite}
  \mathbb{E}[\Delta_V^{\text{sgl}}]
  \;=\; \frac{1}{|\varrho|}\cdot\frac{2\,\Gamma(\ffd/2)}{\sqrt{\pi}\,\Gamma((\ffd-1)/2)\,(\ffd-1)} - 1,
\end{equation}
which is numerically indistinguishable from Eq.~\eqref{eq:theorem-e5}
for all $\ffd \geq 3{,}000$ in our model set.



% ============================================================
% NeurIPS paper checklist
% ============================================================
\newpage
\input{checklist}

\end{document}
